% !TEX root = ../TRENCH_DIFFEQ.tex

\chapter{Linear Higher Order Equations}\label{ch:higher_order}

\noindent
\textsc{In this chapter}
 we  extend  the results obtained
in \zcref{ch:second_order} for linear second order  equations
to linear higher order equations.

\medskip\noindent
\zcref[S,CAP,font=\scshape]{sec:higher_theory} presents a theoretical introduction to
linear higher order equations.


\medskip\noindent
\zcref[S,CAP,font=\scshape]{sec:higher_const_coeff} discusses higher order constant
coefficient
homogeneous equations.


\medskip\noindent
\zcref[S,CAP,font=\scshape]{sec:higher_undetermined} presents the method of
undetermined coefficients for
higher order equations.


\medskip\noindent
\zcref[S,CAP,font=\scshape]{sec:higher_variation_param} extends the
method of  variation of parameters to  higher order
equations.

\newpage

\section{Introduction to Linear Higher Order Equations}
\label{sec:higher_theory}

An $n$th order differential equation is said to be \emph{\color{blue}linear} if it
can be written in the form
\begin{equation}\label{eq:9.1.1}
y^{(n)}+p_1(x)y^{(n-1)}+\cdots+p_n(x)y=f(x).
\end{equation}
We considered equations of this form with $n=1$ in
\zcref{sec:linar_first_order} and with $n=2$ in
\zcref{ch:second_order}. In this chapter $n$ is an arbitrary
positive integer.

In this section we sketch the general theory of linear $n$th order
equations. Since this theory has already been discussed for $n=2$ in
\zcref{sec:homog_linear,sec:nonhomog_linear}, we'll omit
proofs.

For convenience, we consider linear differential equations written
as
\begin{equation}\label{eq:9.1.2}
P_0(x)y^{(n)}+P_1(x)y^{(n-1)}+\cdots+P_n(x)y=F(x),
\end{equation}
which can be rewritten as \eqref{eq:9.1.1} on any interval on
which $P_0$ has no zeros, with $p_1=P_1/P_0$, \dots, $p_n=P_n/P_0$ and
$f=F/P_0$. For simplicity, throughout this chapter we'll
abbreviate the left side of \eqref{eq:9.1.2} by $Ly$; that is,
\[
Ly=P_0y^{(n)}+P_1y^{(n-1)}+\cdots+P_ny.
\]
We say that the equation $Ly=F$ is \emph{\color{blue}normal} on $(a,b)$ if
$P_0$, $P_1$, \dots, $P_n$ and $F$ are continuous on $(a,b)$ and $P_0$ has
no zeros on $(a,b)$. If this is so then $Ly=F$ can be
written as \eqref{eq:9.1.1} with $p_1$, \dots, $p_n$ and $f$ continuous on
$(a,b)$.

The next theorem is analogous to
\zcref{thmtype:5.3.1}.

\begin{theorem}\label{thmtype:9.1.1}\index{uniqueness theorem}
Suppose   $Ly=F$ is normal on $(a,b)$, let $x_0$ be a point in
$(a,b)$,
and let $k_0$, $k_1$, \dots, $k_{n-1}$ be arbitrary real numbers. Then
the initial value problem
\[
Ly=F, \quad  y(x_0)=k_0,\quad y'(x_0)=k_1,\dots,\quad y^{(n-1)}(x_0)=k_{n-1}
\]
has a unique solution on $(a,b)$.
\end{theorem}


\boxit{Homogeneous Equations}

Eqn.~\eqref{eq:9.1.2} is said to be \emph{\color{blue}homogeneous} if $F\equiv0$ and
\emph{\color{blue}nonhomogeneous} otherwise. Since $y\equiv0$ is obviously a
solution of $Ly=0$, we call it the \emph{\color{blue}trivial} solution. Any other
solution is \emph{\color{blue}nontrivial}.



If $y_1$, $y_2$, \dots, $y_n$ are defined on $(a,b)$ and
$c_1$, $c_2$, \dots, $c_n$ are constants, then
\begin{equation}\label{eq:9.1.3}
y=c_1y_1+c_2y_2+\cdots+c_ny_n
\end{equation}
is a \emph{\color{blue}linear combination} of $\{y_1,y_2,\dots,y_n\}$. It's
easy
 to show that if $y_1$, $y_2$, \dots, $y_n$ are solutions of
$Ly=0$ on $(a,b)$, then so is any linear combination of
$\{y_1,y_2,\dots,y_n\}$. (See the proof of \zcref{thmtype:5.1.2}.) We
say that $\{y_1,y_2,\dots,y_n\}$ is a \emph{\color{blue}fundamental set of
solutions of $Ly=0$ on} $(a,b)$ if every solution of $Ly=0$ on
$(a,b)$ can be written as a linear combination of
$\{y_1,y_2,\dots,y_n\}$, as in \eqref{eq:9.1.3}. In this case we say that
\eqref{eq:9.1.3} is the \emph{\color{blue}general solution of $Ly=0$ on
$(a,b)$}.


It can be shown (\zcref{exer:9.1.14,exer:9.1.15}) that if
the equation $Ly=0$ is normal on $(a,b)$ then it has
infinitely many fundamental  sets of solutions  on $(a,b)$.
The next definition will help to identify fundamental sets of
solutions of $Ly=0$.



We say that $\{y_1,y_2,\dots,y_n\}$ is \emph{\color{blue}linearly independent} on
$(a,b)$ if the only constants $c_1$, $c_2$, \dots, $c_n$ such that
\begin{equation}\label{eq:9.1.4}
c_1y_1(x)+c_2y_2(x)+\cdots+c_ny_n(x)=0,\quad a<x<b,
\end{equation}
are $c_1=c_2=\cdots=c_n=0$. If \eqref{eq:9.1.4} holds for some set of
constants $c_1$, $c_2$, \dots, $c_n$ that are not all zero, then
$\{y_1,y_2,\dots,y_n\}$ is \emph{\color{blue}linearly dependent} on $(a,b)$

The next theorem is analogous to
\zcref{thmtype:5.1.3}.

\begin{theorem}\color{blue}\label{thmtype:9.1.2}
If  $Ly=0$ is normal on $(a,b)$, then a set
$\{y_1,y_2,\dots,y_n\}$ of $n$ solutions of $Ly=0$ on $(a,b)$ is a
fundamental set if and only if it's linearly independent on $(a,b)$.
\end{theorem}


\begin{example}\label{example:9.1.1}
The equation
\begin{equation}\label{eq:9.1.5}
x^3y'''-x^2y''-2xy'+6y=0
\end{equation}
is normal and has the solutions $y_1=x^2$, $y_2=x^3$, and $y_3=1/x$ on
$(-\infty,0)$ and $(0,\infty)$.  Show that $\{y_1,y_2,y_3\}$
is linearly independent on $(-\infty, 0)$ and $(0,\infty)$. Then find
the general solution of \eqref{eq:9.1.5} on $(-\infty, 0)$ and
$(0,\infty)$.
\end{example}

\begin{solution}
Suppose
\begin{equation} \label{eq:9.1.6}
c_1x^2+c_2x^3+\frac{c_3}{x}=0
\end{equation}
on $(0,\infty)$. We must show that $c_1=c_2=c_3=0$. Differentiating
\eqref{eq:9.1.6} twice yields  the system
\begin{equation} \label{eq:9.1.7}
\begin{aligned}
c_1x^2+c_2x^3+\dfrac{c_3}{x}&=0\\
2c_1x+3c_2x^2-\dfrac{c_3}{x^2}&=0\\
2c_1+6c_2x+\dfrac{2c_3}{x^3}&=0.
\end{aligned}
\end{equation}
If \eqref{eq:9.1.7} holds for all $x$ in $(0,\infty)$, then it
certainly  holds at $x=1$;  therefore,
\begin{equation} \label{eq:9.1.8}
\begin{aligned}
\phantom{2}c_1+\phantom{3}c_2+\phantom{2}c_3&=0\\
2c_1+3c_2-\phantom{2}c_3&=0\\
2c_1+6c_2+2c_3&=0.
\end{aligned}
\end{equation}
By solving this system directly, you can verify that it has only the
trivial solution $c_1=c_2=c_3=0$;     however, for our purposes it's more
useful to recall from linear algebra that a homogeneous linear system
of $n$ equations in $n$ unknowns has only the trivial solution if its
determinant is nonzero. Since the determinant of \eqref{eq:9.1.8} is
\[
\begin{vmatrix}1&1&1\\2&3&-1\\2&6&2\end{vmatrix}=
\begin{vmatrix}1&0&0\\2&1&-3\\2&4&0\end{vmatrix}=12,
\]
it follows that \eqref{eq:9.1.8} has only the trivial solution, so
$\{y_1,y_2,y_3\}$ is linearly independent on $(0,\infty)$. Now
\zcref{thmtype:9.1.2} implies that
\[
y=c_1x^2+c_2x^3+\frac{c_3}{x}
\]
is the general solution of \eqref{eq:9.1.5} on $(0,\infty)$.
To see that this is also true
 on $(-\infty,0)$,  assume that \eqref{eq:9.1.6} holds on
$(-\infty,0)$. Setting $x=-1$ in \eqref{eq:9.1.7} yields
\begin{align*}
\phantom{-2}c_1-\phantom{3}c_2-\phantom{2}c_3&=0\\
-2c_1+3c_2-\phantom{2}c_3&=0\\
\phantom{-}2c_1-6c_2-2c_3&=0.
\end{align*}
Since the determinant of this system is
\[
\begin{vmatrix}1&-1&-1\\-2&3&-1\\2&-6&-2\end{vmatrix}=
\begin{vmatrix}1&0&0\\-2&1&-3\\2&-4&0\end{vmatrix}=-12,
\]
it follows that $c_1=c_2=c_3=0$; that is, $\{y_1,y_2,y_3\}$
is linearly independent on $(-\infty,0)$.
\end{solution}

\begin{example}\label{example:9.1.2}
The equation
\begin{equation}\label{eq:9.1.9}
y^{(4)}+y'''-7y''-y'+6y=0
\end{equation}
is normal and has the solutions $y_1=e^x$, $y_2=e^{-x}$, $y_3=e^{2x}$,
and $y_4=e^{-3x}$ on $(-\infty,\infty)$. (Verify.) Show that
$\{y_1,y_2,y_3,y_4\}$ is linearly independent on $(-\infty,\infty)$.
Then find the general solution of \eqref{eq:9.1.9}.
\end{example}

\begin{solution}
Suppose   $c_1$, $c_2$, $c_3$, and $c_4$
are constants such that
\begin{equation} \label{eq:9.1.10}
c_1e^x+c_2e^{-x}+c_3e^{2x}+c_4e^{-3x}=0
\end{equation}
for all $x$. We must show that $c_1=c_2=c_3=c_4=0$. Differentiating
\eqref{eq:9.1.10} three times yields the system
\begin{equation} \label{eq:9.1.11}
\begin{aligned}
c_1e^x+c_2e^{-x}+\phantom{2}c_3e^{2x}+\phantom{27}c_4e^{-3x}&=0\\
c_1e^x-c_2e^{-x}+2c_3e^{2x}-\phantom{2}3c_4e^{-3x}&=0\\
c_1e^x+c_2e^{-x}+4c_3e^{2x}+\phantom{2}9c_4e^{-3x}&=0\\
c_1e^x-c_2e^{-x}+8c_3e^{2x}-27c_4e^{-3x}&=0.
\end{aligned}
\end{equation}
If \eqref{eq:9.1.11} holds for all $x$, then it certainly holds for $x=0$.
Therefore
\begin{align*}
c_1+c_2+\phantom{2}c_3+\phantom{27}c_4&=0\\
c_1-c_2+2c_3-\phantom{2}3c_4&=0\\
c_1+c_2+4c_3+\phantom{2}9c_4&=0\\
c_1-c_2+8c_3-27c_4&=0.
\end{align*}
The determinant of this system is
\begin{equation} \label{eq:9.1.12}
\begin{aligned}
\begin{vmatrix}1&1&1&1\\1&-1&2&-3\\1&1&4&9\\1&-1&8&-27\end{vmatrix}
&=\begin{vmatrix}1&1&1&1\\0&-2&1&-4\\0&0&3&8\\0&-2&7&-28\end{vmatrix}
=\begin{vmatrix}-2&1&-4\\0&3&8\\-2&7&-28\end{vmatrix}\\
&=\begin{vmatrix}-2&1&-4\\0&3&8\\0&6&-24\end{vmatrix}
=-2\begin{vmatrix}3&8\\6&-24\end{vmatrix}=240,
\end{aligned}
\end{equation}
so the system has only the trivial solution $c_1=c_2=c_3=c_4=0$.
Now \zcref{thmtype:9.1.2} implies that
\[
y=c_1e^x+c_2e^{-x}+c_3e^{2x}+c_4e^{-3x}
\]
is the general solution of \eqref{eq:9.1.9}.
\end{solution}


\boxit{The Wronskian}\index{Wronskian}

We can use the method used in \zcref{example:9.1.1,example:9.1.2} to test $n$ solutions $\{y_1,y_2,\dots,y_n\}$ of any
$n$th order equation $Ly=0$ for linear independence on an interval
$(a,b)$ on which the equation is normal. Thus, if $c_1$, $c_2$ ,\dots,
$c_n$ are constants such that
\[
c_1y_1+c_2y_2+\cdots+c_ny_n=0,\quad a<x<b,
\]
then differentiating $n-1$ times leads to the $n\times n$
system of equations
\begin{equation} \label{eq:9.1.13}
\begin{aligned}
c_1y_1(x)+c_2y_2(x)+&\cdots+c_ny_n(x)=0\\
c_1y_1'(x)+c_2y_2'(x)+&\cdots+c_ny_n'(x)=0\\
&\vdots\\
c_1y_1^{(n-1)}(x)+c_2y_2^{(n-1)}(x)+&\cdots+c_ny_n^{(n-1)}(x)=0
\end{aligned}
\end{equation}
for $c_1$, $c_2$, \dots, $c_n$. For a fixed $x$,  the determinant of this
system is
\[
W(x)=\begin{vmatrix}
y_1(x)&y_2(x)&\cdots&y_n(x)\\
y'_1(x)&y'_2(x)&\cdots&y_n'(x)\\
\vdots&\vdots&\ddots&\vdots\\
y_1^{(n-1)}(x)&y_2^{(n-1)}(x)&\cdots&y_n^{(n-1)}(x)
\end{vmatrix}.
\]
We call this determinant the
\personHref{Wronski}{Wronskian} of
$\{y_1,y_2,\dots,y_n\}$. If $W(x)\ne0$ for some $x$ in $(a,b)$ then
 the system \eqref{eq:9.1.13} has only the trivial solution
$c_1=c_2=\cdots=c_n=0$, and \zcref{thmtype:9.1.2} implies that
\[
y=c_1y_1+c_2y_2+\cdots+c_ny_n
\]
is the general solution of $Ly=0$ on $(a,b)$.


The next theorem generalizes \zcref{thmtype:5.1.4}.
The proof is sketched in (\zcref[range]{exer:9.1.17,exer:9.1.20}).

\begin{theorem}\label{thmtype:9.1.3}
 Suppose the   homogeneous linear $n$th order equation
\begin{equation}\label{eq:9.1.14}
P_0(x)y^{(n)}+P_1(x)y^{n-1}+\cdots+P_n(x)y=0
\end{equation}
is normal on $(a,b)$, let $y_1$, $y_2$, \dots, $y_n$ be solutions of
\eqref{eq:9.1.14} on $(a,b)$, and let $x_0$ be in $(a,b)$. Then the
Wronskian of $\{y_1,y_2,\dots,y_n\}$ is given by
\begin{equation}  \label{eq:9.1.15}
W(x)=W(x_0)\exp\left\{-\int^x_{x_0}\frac{P_1(t)}{P_0(t)}\,
dt\right\},\quad a<x<b.
\end{equation}
Therefore,  either $W$ has no zeros in  $(a,b)$ or $W\equiv0$
on  $(a,b)$.
\end{theorem}


Formula \eqref{eq:9.1.15} is
\personHref{Abel}{Abel's formula}\index{Abel's formula}.

The next theorem is analogous to
\zcref{thmtype:5.1.6}.

\begin{theorem}\color{blue}\label{thmtype:9.1.4}
Suppose $Ly=0$ is normal on $(a,b)$ and let $y_1$, $y_2$, \dots, $y_n$
be $n$ solutions of $Ly=0$ on $(a,b)$. Then the following statements
are equivalent;  that is, they are either all true or all false:
\begin{alist}
\item % (a)
The general solution of $Ly=0$ on $(a,b)$ is
$y=c_1y_1+c_2y_2+\cdots+c_ny_n$.
\item % (b)
 $\{y_1,y_2,\dots,y_n\}$ is a fundamental set of solutions of
$Ly=0$ on $(a,b)$.
\item % (c)
$\{y_1,y_2,\dots,y_n\}$ is linearly independent on $(a,b)$.
\item % (d)
The Wronskian of $\{y_1,y_2,\dots,y_n\}$ is nonzero at some point in
$(a,b)$.
\item % (e)
The Wronskian of $\{y_1,y_2,\dots,y_n\}$ is nonzero at all points in
$(a,b)$.
\end{alist}
\end{theorem}

\begin{example}\label{example:9.1.3}
In \zcref{example:9.1.1} we saw that the solutions $y_1=x^2$,
$y_2=x^3$, and $y_3=1/x$ of
\[
x^3y'''-x^2y''-2xy'+6y=0
\]
are linearly independent on $(-\infty,0)$ and $(0,\infty)$. Calculate
the Wronskian of $\{y_1,y_2,y_3\}$.
\end{example}

\begin{solution}
If $x\ne0$, then
\[
W(x)=\begin{vmatrix}
x^2&\phantom{2}x^3&1/x\\
2x&3x^2\phantom{^2}&-1/x^2\\
2&6x\phantom{x^2}&2/x^3
\end{vmatrix}
=2x^3\begin{vmatrix}
1&\phantom{3}x&1/x^3\\
2&3x&-1/x^3\\
1&3x&1/x^3
\end{vmatrix},
\]
where we factored $x^2$, $x$, and $2$ out of the first, second, and
third  rows of
$W(x)$, respectively. Adding the second row of the last determinant
to the first and third rows yields
\[
W(x)=2x^3\begin{vmatrix}
3&4x&0\\
2&3x&-1/x^3\\
3&6x&0
\end{vmatrix}
=2x^3\left(\frac{1}{x^3}\right)\begin{vmatrix}3&4x\\3&6x\end{vmatrix}
=12x.
\]
Therefore $W(x)\ne0$ on $(-\infty,0)$ and $(0,\infty)$.
\end{solution}

\begin{example}\label{example:9.1.4}
In
\zcref{example:9.1.2} we saw that the solutions $y_1=e^x$,
$y_2=e^{-x}$,
$y_3=e^{2x}$, and $y_4=e^{-3x}$ of
\[
y^{(4)}+y'''-7y''-y'+6y=0
\]
are linearly independent on every open interval.
Calculate the Wronskian  of $\{y_1,y_2,y_3,y_4\}$.
\end{example}

\begin{solution}
For all $x$,
\[
W(x)=\begin{vmatrix}
e^x&e^{-x}&e^{2x}&e^{-3x}\\
e^x&-e^{-x}&2e^{2x}&-3e^{-3x}\\
e^x&e^{-x}&4e^{2x}&9e^{-3x}\\
e^x&-e^{-x}&8e^{2x}&-27e^{-3x}
\end{vmatrix}.
\]
Factoring the exponential common factor from each row yields
\[
W(x)
=e^{-x}\begin{vmatrix}1&1&1&1\\1&-1&2&-3\\1&1&4&9\\1&-1&8&-27\end{vmatrix}
=240e^{-x},
\]
from \eqref{eq:9.1.12}.
\end{solution}

\begin{remark*}
Under the assumptions of \zcref{thmtype:9.1.4}, it isn't
necessary to obtain a formula for $W(x)$. Just evaluate $W(x)$
at a convenient point in $(a,b)$, as we did in
\zcref{example:9.1.1,example:9.1.2}.
\end{remark*}


\begin{theorem}\label{thmtype:9.1.5}
Suppose $c$ is in $(a,b)$ and $\alpha_{1}$, $\alpha_{2}$, \dots,
are real numbers, not all zero.
Under the assumptions of \zcref{thmtype:10.3.3}, suppose
$y_{1}$ and $y_{2}$ are  solutions  of \eqref{eq:5.1.35}   such that
\begin{equation} \label{eq:9.1.16}
\alpha y_{i}(c)+ y_{i}'(c)+\cdots +y_{i}^{(n-1)}(c)=0,\quad 1\le i\le n.
\end{equation}
Then $\{y_{1},y_{2},\dots y_{n}\}$ isn't  linearly independent on $(a,b)$.
\end{theorem}


\begin{proof}
Since $\alpha_{1}$, $\alpha_{2}$, \dots, $\alpha_{n}$
are not all zero,
\eqref{eq:9.1.14} implies that
\[
\begin{vmatrix}
y_{1}(c)&y_{1}'(c)&\cdots&y_{1}^{(n-1)}(c)\\
y_{2}(c)&y_{2}'(c)&\cdots&y_{2}^{(n-1)}(c)\\
\vdots&\vdots&\ddots&\vdots\\
y_{n}(c)&y_{n}'(c)&\cdots&y_{n}^{(n-1)}(c)\\
\end{vmatrix}=0,
\]
so
\[
\begin{vmatrix}
y_{1}(c)&y_{2}(c)&\cdots& y_{n}(c)\\
y_{1}'(c)&y_{2}'(c)&\cdots& y_{n}'(c)\\
\vdots&\vdots&\ddots&\vdots\\
y_{1}^{(n-1)}(c)&y_{2}^{(n-1)}(c)(c)&\cdots& y_{n}^{(n-1)}(c)(c)\\
\end{vmatrix}=0
\]
and \zcref{thmtype:9.1.4} implies the stated conclusion.
\end{proof}


\boxit{General Solution of a Nonhomogeneous Equation}

The next theorem is analogous to
\zcref{thmtype:5.3.2}.
It shows how to find the general solution of $Ly=F$ if we know a particular
solution of $Ly=F$ and a fundamental set of solutions of the
\indexterm{complementary equation} $Ly=0$.

\begin{theorem}\color{blue}\label{thmtype:9.1.6}
Suppose $Ly=F$ is normal on $(a,b)$. Let $y_p$ be a particular
solution of $Ly=F$ on $(a,b)$, and let $\{y_1,y_2,\dots,y_n\}$ be a
fundamental set of solutions of the complementary equation $Ly=0$ on
$(a,b)$. Then $y$ is a solution of $Ly=F$ on $(a,b)$ if and only if
\[
y=y_p+c_1y_1+c_2y_2+\cdots+c_ny_n,
\]
where $c_1,c_2,\dots,c_n$  are constants.
\end{theorem}

The next theorem is analogous to \zcref{thmtype:5.3.2}.

\begin{theorem}[The Principle of Superposition]\label{thmtype:9.1.7}
\index{principle of superposition}\index{superposition!principle of}
Suppose for each $i=1$, $2$, \dots, $r$, the function $y_{p_i}$ is a
particular solution of $Ly=F_i$ on  $(a,b)$. Then
\[
y_p=y_{p_1}+y_{p_2}+\cdots+y_{p_r}
\]
is a particular  solution of
\[
Ly=F_1(x)+F_2(x)+\cdots+F_r(x)
\]
on $(a,b)$.
\end{theorem}

We'll apply  \zcref{thmtype:9.1.6,thmtype:9.1.7}
throughout the rest of this chapter.

\exercises

\begin{exerciselist}

\item\label{exer:9.1.1}Verify that the given function is the solution of the initial
value problem.

\begin{alist}
\item $ x^3y'''-3x^2y''+6xy'-6y=-\dfrac{24}{x}, \quad  y(-1)=0$,
$y'(-1)=0$, \quad  $y''(-1)=0$; $y=-6x-8x^2-3x^3+\dfrac{1}{x}$
\item $y'''-\dfrac{1}{x}y''-y'+\dfrac{1}{x}y=\dfrac{x^2-4}{x^4}$, \quad
$y(1)=\dfrac{3}{2}$,\quad $y'(1)=\dfrac{1}{2}$, $ y''(1)=1$;
$y=x+\dfrac{1}{2x}$
\item $xy'''-y''-xy'+y=x^2, \quad  y(1)=2,\quad y'(1)=5,\quad y''(1)=-1$;
$y=-x^2-2+2e^{(x-1)}-e^{-(x-1)}+4x$
\item $4x^3y'''+4x^2y''-5xy'+2y=30x^2, \quad
y(1)=5,\quad y'(1)=\dfrac{17}{2}$;
$y''(1)=\dfrac{63}{4}$;\quad $y=2x^2\ln x-x^{1/2}+2x^{-1/2}+4x^2$
\item $x^4y^{(4)}-4x^3y'''+12x^2y''-24xy'+24y=6x^4, \quad  y(1)=-2$,
$y'(1)=-9, \quad  y''(1)=-27,\quad y'''(1)=-52$;
$y=x^4\ln x+x-2x^2+3x^3-4x^4$
\item $xy^{(4)}-y'''-4xy''+4y'=96x^2, \quad  y(1)=-5,\quad y'(1)=-24$
$y''(1)=-36; \quad   y'''(1)=-48;\quad y=9-12x+6x^2-8x^3$
\end{alist}

\item\label{exer:9.1.2}Solve the initial value problem
\[
x^3y'''-x^2 y''-2xy'+6y=0,
\quad   y(-1)=-4, \quad   y'(-1)=-14,\quad   y''(-1)=-20.
\]
\hint{See \zcref{example:9.1.1}.}


\item\label{exer:9.1.3}Solve the initial value problem
\[
y^{(4)}+y'''-7y''-y'+6y=0,
\quad   y(0)=5,\quad   y'(0)=-6,\quad   y''(0)=10,\quad   y'''(0)-36.
\]
\hint{See \zcref{example:9.1.2}.}


\item\label{exer:9.1.4}Find solutions $y_1$, $y_2$, \dots, $y_n$ of the equation $y^{(n)}=0$
that satisfy the initial conditions
\[
y_i^{(j)}(x_0)=\begin{cases}
0,&j\ne i-1,\\
1,&j=i-1,\end{cases}\quad 1\le i\le n.
\]

\item\label{exer:9.1.5}\begin{alist}
\item %(a)
Verify that the function
\[
y=c_1x^3+c_2x^2+\frac{c_3}{x}
\]
satisfies
\[
x^3 y'''-x^2y''-2xy'+6y=0
\tag{A}
\]
if $c_1$, $c_2$, and $c_3$ are constants.

\item %(b)
Use (a) to find solutions $y_1$, $y_2$, and $y_3$ of
(A) such that
\begin{align*}
y_1(1)&=1,\quad y_1'(1)&=0,\quad  y_1''(1)&=0  \\
y_2(1)&=0,\quad  y_2'(1)&=1,\quad  y_2''(1)&=0  \\
y_3(1)&=0,\quad  y_3'(1)&=0,\quad  y_3''(1)&=1.
\end{align*}

\item %(c)
Use (b) to find the solution of  (A) such that
\[
y(1)=k_0,\quad y'(1)=k_1,\quad   y''(1)=k_2.
\]
\end{alist}

\item\label{exer:9.1.6}Verify that the given functions are solutions of the given equation, and
 show that they  form a fundamental set of solutions of the
equation on any interval on which the equation is normal.

\begin{alist}
\item $y'''+y''-y'-y=0;  \quad\{e^x,\,e^{-x},\,xe^{-x}\}$
\item $y'''-3y''+7y'-5y=0;  \quad\{e^x,\,e^x\cos2x,\,e^x\sin2x\}$.
\item $xy'''-y''-xy'+y=0;  \quad \{e^x,\,e^{-x},\,x\}$
\item $x^2y'''+2xy''-(x^2+2)y=0;  \quad\{e^x/ x,\,e^{-x}/x,\,1\}$
\item $(x^2-2x+2)y'''-x^2y''+2xy'-2y=0;  \quad \{x,\,x^2,\,e^x\}$
\item $(2x-1)y^{(4)}-4xy'''+(5-2x)y''+4xy'-4y=0;
\quad\{x,\,e^x,\,e^{-x},e^{2x}\}$
\item $xy^{(4)}-y'''-4xy'+4y'=0;  \quad\{1,x^2,\,e^{2x},\,e^{-2x}\}$
\end{alist}

\item\label{exer:9.1.7}Find the Wronskian $W$ of a set of three solutions
of \[ y'''+2xy''+e^xy'-y=0, \] given that $W(0)=2$.

\item\label{exer:9.1.8}Find the Wronskian $W$ of a set of four solutions of
\[
y^{(4)}+(\tan x)y'''+x^2y''+2xy=0,
\]
given that $W(\pi/4)=K$.

\item\label{exer:9.1.9}\begin{alist}
\item %(a)
Evaluate the Wronskian $W$ $\{e^x,\,xe^x,\,
x^2e^x\}$. Evaluate $W(0)$.

\item %(b)
Verify that $y_1$, $y_2$, and $y_3$ satisfy
\[
y'''-3y''+3y'-y=0.
\tag{A}
\]

\item %(c)
Use $W(0)$ from (a) and Abel's formula to calculate $W(x)$.

\item %(d)
What is the general solution of  (A)?
\end{alist}

\item\label{exer:9.1.10}Compute the Wronskian of the given set of functions.

\begin{alist*}[columns=2]
\item $\{1,\,e^x,\,e^{-x}\}$
\item $\{e^x,\, e^x\sin x,\,e^x\cos x\}$
\item $\{2,\,x+1,\,x^2+2\}$
\item $\{x,\,x\ln x,\,1/x\}$
\item $\{1,\,x,\,\dfrac{x^2}{2!},\,\dfrac{x^3}{3!}\,,\cdots,\,\dfrac{x^n}{n!}\}$
\item $\{e^x,\,e^{-x},\,x\}$
\item $\{e^x/x,\,e^{-x}/x,\,1\}$
\item $\{x,\,x^2,\,e^x\}$
\item $\{x,\,x^3,\,1/x,\,1/x^2\}$
\item $\{e^x,\,e^{-x},\,x,\,e^{2x}\}$
\item $\{e^{2x},\,e^{-2x},\,1,\,x^2\}$
\end{alist*}

\item\label{exer:9.1.11}Suppose $Ly=0$ is normal on $(a,b)$ and $x_0$ is in $(a,b)$. Use
\zcref{thmtype:9.1.1} to show that $y\equiv0$ is the only solution
of the initial value problem
\[
Ly=0, \quad  y(x_0)=0,\quad y'(x_0)=0,\dots, y^{(n-1)}(x_0)=0,
\]
on $(a,b)$.

\item\label{exer:9.1.12}Prove:  If $y_1$, $y_2$, \dots, $y_n$ are solutions of $Ly=0$ and the
functions
\[
z_i=\sum^n_{j=1}a_{ij}y_j,\quad 1\le i\le n,
\]
form a fundamental set of solutions of $Ly=0$, then so do $y_1$, $y_2$,
\dots, $y_n$.

\item\label{exer:9.1.13}Prove:  If
\[
y=c_1y_1+c_2y_2+\cdots+c_ky_k+y_p
\]
is a solution of a linear equation $Ly=F$ for every choice of the
constants $c_1$, $c_2$ ,\dots, $c_k$, then $Ly_i=0$ for $1\le i\le k$.

\item\label{exer:9.1.14}Suppose $Ly=0$ is normal on $(a,b)$ and let $x_0$ be in $(a,b)$.
For $1\le i\le n$, let $y_i$ be the solution of the initial value
problem
\[
Ly_i=0, \quad  y_i^{(j)} (x_0)=
\begin{cases}
0,& j\ne i-1,\\
1,&j=i-1,
\end{cases}
1\le i\le n,
\]
where $x_0$ is an arbitrary point in $(a,b)$.  Show that any solution
of $Ly=0$ on $(a, b)$, can be written as
\[
y=c_1y_1+c_2y_2+\cdots+c_ny_n,
\]
with $c_j=y^{(j-1)}(x_0)$.

\item\label{exer:9.1.15}Suppose $\{y_1, y_2,\dots, y_n\}$ is a fundamental set of
solutions of
\[
P_0(x)y^{(n)}+P_1(x)y^{(n-1)}+\cdots+P_n(x)y=0
\]
on $(a,b)$, and let
\begin{align*}
z_1&=a_{11}y_1+a_{12}y_2+\cdots+a_{1n}y_n\\
z_2&=a_{21}y_1+a_{22}y_2+\cdots+a_{2n}y_n\\
&\vdots\phantom{_1y_1+a}\vdots\phantom{_2y_2+\cdots+a}\vdots\\
z_n&=a_{n1}y_1+a_{n2}y_2+\cdots+a_{nn}y_n,
\end{align*}
where the $\{a_{ij}\}$ are constants. Show that $\{z_1, z_2,\dots,
z_n\}$ is a fundamental set of solutions  of (A) if and only if the
determinant
\[
\begin{vmatrix}
a_{11}&a_{12}&\cdots&a_{1n}\\
a_{21}&a_{22}&\cdots&a_{2n}\\
\vdots&\vdots&\ddots&\vdots\\
a_{n1}&a_{n2}&\cdots&a_{nn}
\end{vmatrix}
\]
is nonzero. \hint{The determinant of a product of  $n\times
n$ matrices equals the product of the determinants.}

\item\label{exer:9.1.16}Show that $\{y_1,y_2,\dots,y_n\}$ is linearly dependent on $(a,b)$ if
and only if at least one of the functions $y_1$, $y_2$, \dots, $y_n$ can be
written as a linear combination of the others on $(a,b)$.

\exercisetext{Take the following as a hint in
\zcref[range]{exer:9.1.17,exer:9.1.19}:\\
\color{blue}
 By the definition of determinant,
\[
\begin{vmatrix}
a_{11}&a_{12}&\cdots&a_{1n}\\
a_{21}&a_{22}&\cdots&a_{2n}\\
\vdots&\vdots&\ddots&\vdots\\
a_{n1}&a_{n2}&\cdots&a_{nn}
\end{vmatrix}
=\sum\pm a_{1i_1}a_{2i_2},\dots,a_{ni_n},
\]
where the sum is over all permutations $(i_1,i_2,\dots,i_n)$ of
$(1,2,\dots, n)$ and the choice of $+$ or $-$ in each term depends
only on the permutation associated with that term.}

\item\label{exer:9.1.17}Prove:  If
\[
A(u_1,u_2,\dots,u_n)=
\begin{vmatrix}
a_{11}&a_{12}&\cdots&a_{1n}\\
a_{21}&a_{22}&\cdots&a_{2n}\\
\vdots&\vdots&\ddots&\vdots\\
a_{n-1,1}&a_{n-1,2}&\cdots&a_{n-1,n}\\
u_1&u_2&\cdots&u_n
\end{vmatrix},
\]
then
\[
A(u_1+v_1, u_2+v_2,\dots,
  u_n+v_n)=A(u_1,u_2,\dots,u_n)+A(v_1,v_2,\dots,
  v_n).
\]

\item\label{exer:9.1.18}Let
\[
F=\begin{vmatrix}
f_{11}&f_{12}&\cdots&f_{1n}\\
f_{21}&f_{22}&\cdots&f_{2n}\\
\vdots&\vdots&\ddots&\vdots\\
f_{n1}&f_{n2}&\cdots&f_{nn}
\end{vmatrix},
\]
where $f_{ij}\; (1\le i,\;   j\le n)$ is differentiable.  Show that
\[
F'=F_1+F_2+\cdots+F_n,
\]
where $F_i$ is the determinant obtained by differentiating the
$i$th row of $F$.


\item\label{exer:9.1.19}Use \zcref{exer:9.1.18}  to show that if $W$ is the
Wronskian of the
$n$-times differentiable functions $y_1$, $y_2$, \dots, $y_n$, then
\[
W'=
\begin{vmatrix}
y_1&y_2&\cdots&y_n\\
y'_1&y'_2&\cdots&y'_n\\
\vdots&\vdots&\ddots&\vdots\\
y_1^{(n-2)}&y_2^{(n-2)}&\cdots&y_n^{(n-2)}\\
y_1^{(n)}&y_2^{(n)}&\cdots&y_n^{(n)}
\end{vmatrix}.
\]

\item\label{exer:9.1.20}Use \zcref{exer:9.1.17,exer:9.1.19} to show that if $W$ is
the Wronskian
of  solutions $\{y_1,y_2,\dots,y_n\}$ of the normal equation
\[
P_0(x)y^{(n)}+P_1(x)y^{(n-1)}+\cdots+P_n(x)y=0,
\tag{A}
\]
then $W'=-P_1W/P_0$.
Derive Abel's formula (Eqn.~\eqref{eq:9.1.15}) from this.
\hint{Use (A) to write $y^{(n)}$ in terms of
$y,y',\dots,y^{(n-1)}$.}

\item\label{exer:9.1.21}Prove \zcref{thmtype:9.1.6}.

\item\label{exer:9.1.22}Prove \zcref{thmtype:9.1.7}.



\item\label{exer:9.1.23}Show that if the Wronskian of the $n$-times continuously
differentiable functions $\{y_1,y_2,\dots,y_n\}$ has no zeros in
$(a,b)$, then the differential equation obtained by expanding the
determinant
\[
\begin{vmatrix}
y&y_1&y_2&\cdots&y_n\\
y'&y'_1&y'_2&\cdots&y'_n\\
\vdots&\vdots&\vdots&\ddots& \vdots\\
y^{(n)}&y_{1}^{(n)}&y_2^{(n)}&\cdots&y_n^{(n)}
\end{vmatrix}=0,
\]
in cofactors of its first column is normal and has
$\{y_1,y_2,\dots,y_n\}$ as a fundamental set of solutions on $(a,b)$.

\item\label{exer:9.1.24}Use the method suggested by \zcref{exer:9.1.23}  to find a linear
homogeneous equation such that the given set of functions
is a fundamental set of solutions on intervals on which the Wronskian
of the set has no zeros.
\begin{alist*}[columns=2,beginpenalty={5000}]
\item $\{x,\,x^2-1,\,x^2+1\}$
\item $\{e^x,\,e^{-x},\,x\}$
\item $\{e^x,\,xe^{-x},\,1\}$
\item $\{x,\,x^2,\,e^x\}$
\item $\{x,\,x^2,\,1/x\}$
\item $\{x+1,\,e^x,\,e^{3x}\}$
\item $\{x,\,x^3,\,1/x,\,1/x^2\}$
\item $\{x,\,x\ln x,\,1/x,\,x^2\}$
\item $\{e^x,\,e^{-x},\,x,\,e^{2x}\}$
\item $\{e^{2x},\,e^{-2x},\,1,\,x^2\}$
\end{alist*}

\end{exerciselist}


\section{Higher Order Constant Coefficient Homogeneous Equations}
\label{sec:higher_const_coeff}

If $a_0$, $a_1$, \dots, $a_n$ are constants and $a_0\ne0$, then
\[
a_0y^{(n)}+a_1y^{(n-1)}+\cdots+a_ny=F(x)
\]
is said to be a \emph{\color{blue}constant coefficient equation}.
In this section we consider the homogeneous constant coefficient
equation
\begin{equation} \label{eq:9.2.1}
a_0y^{(n)}+a_1y^{(n-1)}+\cdots+a_ny=0.
\end{equation}
Since \eqref{eq:9.2.1} is normal on $(-\infty,\infty)$, the theorems in
\zcref{sec:higher_theory} all apply with $(a,b)=(-\infty,\infty)$.

As in \zcref{sec:const_coeff_homog}, we call
\begin{equation} \label{eq:9.2.2}
p(r)=a_0r^n+a_1r^{n-1}+\cdots+a_n
\end{equation}
the \indexterm{characteristic polynomial}\index{polynomial!characteristic} of \eqref{eq:9.2.1}. We saw in
\zcref{sec:const_coeff_homog} that when $n=2$ the solutions of \eqref{eq:9.2.1} are
determined by the zeros of the characteristic polynomial. This is also
true when $n>2$, but the situation is more complicated in this case.
Consequently, we take a different approach here than in
\zcref{sec:const_coeff_homog}.

If $k$ is a positive integer, let $D^k$ stand
for the $k$-th derivative operator;   that is
\[
D^ky=y^{(k)}.
\]
If
\[
q(r)=b_0r^m+b_1r^{m-1}+\cdots+b_m
\]
is an arbitrary polynomial,  define the operator
\[
q(D)=b_0D^m+b_1D^{m-1}+\cdots+b_m
\]
such that
\[
q(D)y=(b_0D^m+b_1D^{m-1}+\cdots+b_m)y=b_0y^{(m)}+b_1y^{(m-1)}+\cdots+
b_my
\]
whenever $y$ is a function with $m$ derivatives. We call $q(D)$
a \indexterm{polynomial operator}.

With $p$ as in \eqref{eq:9.2.2},
\[
p(D)=a_0D^n+a_1D^{n-1}+\cdots+a_n,
\]
so \eqref{eq:9.2.1} can be written as $p(D)y=0$.  If $r$ is a constant
then
\begin{align*}
p(D)e^{rx}&=\left(a_0D^ne^{rx}+a_1D^{n-1}e^{rx}+\cdots+a_ne^{rx}\right)\\
&=(a_0r^n+a_1r^{n-1}+\cdots+a_n)e^{rx};
\end{align*}
that is
\[
p(D)(e^{rx})=p(r)e^{rx}.
\]
This shows that $y=e^{rx}$ is a solution of \eqref{eq:9.2.1}
if $p(r)=0$.
In the simplest case, where $p$ has $n$ distinct real zeros
$r_1$, $r_2$,\dots, $r_n$, this argument yields $n$ solutions
\[
y_1=e^{r_1x},\quad y_2=e^{r_2x},\dots,\quad y_n=e^{r_nx}.
\]
It can be shown (\zcref{exer:9.2.39}) that the Wronskian of
$\{e^{r_1x},e^{r_2x},\dots,e^{r_nx}\}$ is nonzero if
$r_1$, $r_2$, \dots, $r_n$ are distinct; hence,
$\{e^{r_1x},e^{r_2x},\dots,e^{r_nx}\}$
is a fundamental set of solutions of $p(D)y=0$ in this case.

\begin{example}\label{example:9.2.1}
\begin{alist}
\item %(a)
Find the general solution of
\begin{equation} \label{eq:9.2.3}
y'''-6y''+11y'-6y=0.
\end{equation}
\item %(b)
Solve the initial value problem
\begin{equation} \label{eq:9.2.4}
y'''-6y''+11y'-6y=0, \quad  y(0)=4,\quad y'(0)=5,\quad y''(0)=9.
\end{equation}
\end{alist}
\end{example}

\begin{alistSolution}
\item The characteristic polynomial of
\eqref{eq:9.2.3} is
\[
p(r)=r^3-6r^2+11r-6=(r-1)(r-2)(r-3).
\]
Therefore $\{e^x,e^{2x},e^{3x}\}$ is a set of solutions of
\eqref{eq:9.2.3}. It is a fundamental set, since its Wronskian is
\[
W(x)=\begin{vmatrix}
e^x&e^{2x}&e^{3x}\\
e^x&2e^{2x}&3e^{3x}\\
e^x&4e^{2x}&9e^{3x}
\end{vmatrix}=
e^{6x}\begin{vmatrix}1&1&1\\1&2&3\\1&4&9\end{vmatrix}=2e^{6x}\ne0.
\]
Therefore the general solution of \eqref{eq:9.2.3} is
\begin{equation} \label{eq:9.2.5}
y=c_1e^{x}+c_2e^{2x}+c_3e^{3x}.
\end{equation}

\item We must determine $c_1$, $c_2$ and $c_3$ in
\eqref{eq:9.2.5}
so that $y$ satisfies the initial conditions in \eqref{eq:9.2.4}.
Differentiating \eqref{eq:9.2.5} twice yields
\begin{equation} \label{eq:9.2.6}
\begin{aligned}
y'&=c_1e^{x}+2c_2e^{2x}+3c_3e^{3x}\\
y''&=c_1e^{x}+4c_2e^{2x}+9c_3e^{3x}.
\end{aligned}
\end{equation}
Setting $x=0$ in \eqref{eq:9.2.5} and \eqref{eq:9.2.6}  and imposing the
initial conditions yields
\begin{align*}
c_1+\phantom{2}c_2+\phantom{3}c_3&=4\\
c_1+2c_2+3c_3&=5\\
c_1+4c_2+9c_3&=9.
\end{align*}
The solution of this system is $c_1=4$, $c_2=-1$, $c_3=1$. Therefore the
solution of \eqref{eq:9.2.4} is
 \[
y=4e^x-e^{2x}+e^{3x}
\]
(\zcref{figure:9.2.1}).
\end{alistSolution}

\begin{figure}
  \centering\color{blue}
  \includegraphics[alt={A graph of the solution for 0<x<2 and 0<y<700 that begins at (0,4), increases slowly at first, and increases rapidly as x increases.},
  width=5.67in,height=3.3in,keepaspectratio]{fig090201}
\caption{$y=4e^x-e^{2x}+e^{3x}$}
  \label{figure:9.2.1}
\end{figure}

Now we consider the case where the characteristic polynomial
\eqref{eq:9.2.2} does not have $n$ distinct real zeros. For this purpose
it is useful to define what we mean by a factorization of a polynomial
operator. We begin with an example.

\begin{example}\label{example:9.2.2}
Consider the polynomial
\[
p(r)=r^3-r^2+r-1
\]
and the associated polynomial operator
\[
p(D)=D^3-D^2+D-1.
\]
Since $p(r)$ can be factored as
\[
p(r)=(r-1)(r^2+1)=(r^2+1)(r-1),
\]
it's reasonable to expect that p(D) can be factored as
\begin{equation} \label{eq:9.2.7}
p(D)=(D-1)(D^2+1)=(D^2+1)(D-1).
\end{equation}
However, before we can make this assertion we must \emph{\color{blue}define}
 what we mean by saying that two operators are equal, and what we mean
by the products of operators in \eqref{eq:9.2.7}. We say that two
operators are equal if they apply to the same functions and always
produce the same result. The definitions of the products in
\eqref{eq:9.2.7} is this: if $y$ is any three-times differentiable
function then
\begin{alist}
\item % (a)
$(D-1)(D^2+1)y$
is the function obtained by first applying $D^2+1$ to $y$ and then
applying $D-1$ to the resulting function
\item % (b)
$(D^2+1)(D-1)y$
is the function obtained by first applying $D-1$ to $y$ and then
applying $D^2+1$ to the resulting function.
\end{alist}
From (a),
\begin{equation} \label{eq:9.2.8}
\begin{aligned}
(D-1)(D^2+1)y
&=(D-1)[(D^2+1)y]\\
&=(D-1)(y''+y)=D(y''+y)-(y''+y)\\
&=(y'''+y')-(y''+y)\\
&=y'''-y''+y'-y=(D^3-D^2+D-1)y.
\end{aligned}
\end{equation}
This implies that
\[
(D-1)(D^2+1)=(D^3-D^2+D-1).
\]
From (b),
\begin{gather} \label{eq:9.2.9}
\begin{aligned}
(D^2+1)(D-1)y
&=(D^2+1)[(D-1)y]\\
&=(D^2+1)(y'-y)=D^2(y'-y)+(y'-y)\\
&=(y'''-y'')+(y'-y)\\
&=y'''-y''+y'-y=(D^3-D^2+D-1)y,
\end{aligned}
\\
(D^2+1)(D-1)=(D^3-D^2+D-1),\notag
\end{gather}
which completes the justification of \eqref{eq:9.2.7}.
\end{example}

\begin{example}\label{example:9.2.3}
Use the result of \zcref{example:9.2.2} to find the general
solution of
\begin{equation} \label{eq:9.2.10}
y'''-y''+y'-y=0.
\end{equation}
\end{example}

\begin{solution}
From \eqref{eq:9.2.8}, we can rewrite \eqref{eq:9.2.10} as
\[
(D-1)(D^2+1)y=0,
\]
which implies that any solution of $(D^2+1)y=0$ is a solution of
\eqref{eq:9.2.10}. Therefore $y_1=\cos x$ and $y_2=\sin x$ are solutions
of \eqref{eq:9.2.10}.

From \eqref{eq:9.2.9}, we can rewrite \eqref{eq:9.2.10} as
\[
(D^2+1)(D-1)y=0,
\]
which implies that any solution of $(D-1)y=0$ is a solution of
\eqref{eq:9.2.10}. Therefore $y_3=e^x$ is solution of \eqref{eq:9.2.10}.

The Wronskian of $\{e^x,\cos x,\sin x\}$ is
\[
W(x)=\begin{vmatrix}
\cos x&\sin x&e^x\\
-\sin x&\cos x&e^x\\
-\cos x&-\sin x&e^x
\end{vmatrix}.
\]
Since
\[
W(0)=\begin{vmatrix}1&0&1\\0&1&1\\-1&0&1\end{vmatrix}=2,
\]
$\{\cos x,\sin x,e^x\}$ is linearly independent
and
\[
y=c_1\cos x+c_2\sin x+c_3e^x
\]
is the general solution of \eqref{eq:9.2.10}.
\end{solution}

\begin{example}\label{example:9.2.4}
Find the general solution of
\begin{equation} \label{eq:9.2.11}
y^{(4)}-16y=0.
\end{equation}
\end{example}

\begin{solution}
The characteristic polynomial of
\eqref{eq:9.2.11} is
\[
p(r)=r^4-16=(r^2-4)(r^2+4)=(r-2)(r+2)(r^2+4).
\]
By arguments similar to those used in \zcref{example:9.2.2,example:9.2.3}, it can be shown that \eqref{eq:9.2.11} can be written
as
\[
(D^2+4)(D+2)(D-2)y=0
\]
or
\[
(D^2+4)(D-2)(D+2)y=0
\]
or
\[
(D-2)(D+2)(D^2+4)y=0.
\]
Therefore $y$  is a solution of \eqref{eq:9.2.11} if it's a solution
of any  of the three  equations
\[
 (D-2)y=0,\quad (D+2)y=0, \quad(D^2+4)y=0.
\]
Hence, $\{e^{2x},e^{-2x},\cos2x,\sin2x\}$ is a set of solutions of
\eqref{eq:9.2.11}. The Wronskian of this set is
\[
W(x)=\begin{vmatrix}
e^{2x}&e^{-2x}&\cos2x&\sin2x\\
2e^{2x}&-2e^{-2x}&-2\sin2x&2\cos2x\\
4e^{2x}&4e^{-2x}&-4\cos2x&-4\sin2x\\
8e^{2x}&-8e^{-2x}&8\sin2x&-8\cos2x\\
\end{vmatrix}.
\]
Since
\[
W(0)=\begin{vmatrix}
1&1&1&0\\
2&-2&0&2\\
4&4&-4&0\\
8&-8&0&-8\\
\end{vmatrix}=-512,
\]
 $\{e^{2x},e^{-2x},\cos2x,\sin2x\}$ is linearly independent,
and
\[
y_1=c_1e^{2x}+c_2e^{-2x}+c_3\cos2x+c_4\sin2x
\]
is   the general solution of \eqref{eq:9.2.11}.
\end{solution}

It is known from algebra that every polynomial
\[
p(r)=a_0r^n+a_1r^{n-1}+\cdots+a_n
\]
with real coefficients can be factored as
\[
p(r)=a_0p_1(r)p_2(r)\cdots p_k(r),
\]
where no pair of the polynomials $p_1$, $p_2$, \dots, $p_k$  has a commom
factor and each  is either of the form
\begin{equation} \label{eq:9.2.12}
p_j(r)=(r-r_j)^{m_j},
\end{equation}
where $r_j$ is real and $m_j$ is a positive integer, or
\begin{equation} \label{eq:9.2.13}
p_j(r)=\left[(r-\lambda_j)^2+\omega_j^2\right]^{m_j},
\end{equation}
where $\lambda_j$ and $\omega_j$ are real, $\omega_j\ne0$, and $m_j$
is a positive integer. If \eqref{eq:9.2.12} holds then $r_j$ is a real
zero of $p$, while if \eqref{eq:9.2.13} holds then $\lambda+i\omega$
and $\lambda-i\omega$ are complex conjugate zeros of $p$. In either
case, $m_j$ is the \indexterm{multiplicity} of the zero(s).


By arguments similar to those used in our examples, it can be shown
that
\begin{equation} \label{eq:9.2.14}
p(D)=a_0p_1(D)p_2(D)\cdots p_k(D)
\end{equation}
and that the order of the factors on the right can be chosen
arbitrarily. Therefore, if $p_j(D)y=0$ for some $j$ then
$p(D)y=0$. To see this, we simply rewrite \eqref{eq:9.2.14} so that
$p_j(D)$ is applied first. Therefore the problem of finding solutions
of $p(D)y=0$ with $p$ as in \eqref{eq:9.2.14} reduces to finding solutions
of each of these equations
\[
p_j(D)y=0,\quad 1\le j\le k,
\]
where $p_j$ is a power of a first degree term or of an irreducible
quadratic. To find a fundamental set of solutions
$\{y_1,y_2,\dots,y_n\}$ of $p(D)y=0$, we find fundamental set of
solutions of each of the equations and take $\{y_1,y_2,\dots,y_n\}$ to
be the set of all functions in these separate fundamental sets. In
\zcref{exer:9.2.40} we sketch the proof that
$\{y_1,y_2,\dots,y_n\}$ is linearly independent, and therefore a
fundamental set of solutions of $p(D)y=0$.

To apply this procedure to general homogeneous constant coefficient
equations, we must be able to find fundamental sets of solutions of
equations of the form
\[
(D-a)^my=0
\]
and
\[
\left[(D-\lambda)^2+\omega^2\right]^my=0,
\]
where $m$ is an arbitrary positive integer. The next two theorems show
how to do this.

\begin{theorem}\color{blue} \label{thmtype:9.2.1}
If $m$ is a positive integer, then
\begin{equation} \label{eq:9.2.15}
\{e^{ax}, xe^{ax},\dots, x^{m-1}e^{ax}\}
\end{equation}
is a fundamental set of solutions of
\begin{equation} \label{eq:9.2.16}
(D-a)^my=0.
\end{equation}
\end{theorem}

\begin{proof}
We'll show that if
\[
f(x)=c_1+c_2x+\cdots+c_mx^{m-1}
\]
is an arbitrary polynomial of degree $\le m-1$, then $y=e^{ax}f$
is a solution of \eqref{eq:9.2.16}. First note that if $g$ is any
differentiable function then
\[
(D-a)e^{ax}g=De^{ax}g-ae^{ax}g=ae^{ax}g+e^{ax}g'-ae^{ax}g,
\]
so
\begin{equation} \label{eq:9.2.17}
(D-a)e^{ax}g=e^{ax}g'.
\end{equation}
Therefore
\begin{align*}
(D-a)e^{ax}f
&=e^{ax}f'&&\text{(from \eqref{eq:9.2.17} with $g=f$)}\\
(D-a)^2e^{ax}f
&=(D-a)e^{ax}f'=e^{ax}f''&&\text{(from \eqref{eq:9.2.17} with $g=f'$)}\\
(D-a)^3e^{ax}f
&=(D-a)e^{ax}f''=e^{ax}f'''&&\text{(from \eqref{eq:9.2.17} with $g=f''$)}\\
&\vdots\\
(D-a)^me^{ax}f
&=(D-a)e^{ax}f^{(m-1)}=e^{ax}f^{(m)}
&&\text{(from \eqref{eq:9.2.17} with $g=f^{(m-1)}$)}.
\end{align*}
Since $f^{(m)}=0$, the last equation implies that $y=e^{ax}f$ is a
solution of \eqref{eq:9.2.16} if $f$ is any polynomial of degree $\le
m-1$. In particular, each function in \eqref{eq:9.2.15} is a solution of
\eqref{eq:9.2.16}. To see that \eqref{eq:9.2.15} is linearly independent (and
therefore a fundamental set of solutions of \eqref{eq:9.2.16}), note that
if
\[
c_1e^{ax}+c_2xe^{ax}+c\dots+c_{m-1}x^{m-1}e^{ax}=0
\]
for all $x$ in some interval $(a,b)$, then
\[
c_1+c_2x+c\dots+c_{m-1}x^{m-1}=0
\]
for all $x$ in $(a,b)$.  However, we know from algebra that if this
polynomial has more than $m-1$ zeros then $c_1=c_2=\cdots=c_n=0$.
\end{proof}

\begin{example}\label{example:9.2.5}
Find the general solution of
\begin{equation} \label{eq:9.2.18}
y'''+3y''+3y'+y=0.
\end{equation}
\end{example}

\begin{solution}
The characteristic polynomial of \eqref{eq:9.2.18} is
\[
p(r)=r^3+3r^2+3r+1=(r+1)^3.
\]
Therefore \eqref{eq:9.2.18} can be written as
\[
(D+1)^3y=0,
\]
so \zcref{thmtype:9.2.1} implies that the general solution of
\eqref{eq:9.2.18} is
\[
y=e^{-x}(c_1+c_2x+c_3x^2).\qedhere
\]
\end{solution}

The proof of the next theorem is sketched in
\zcref{exer:9.2.41}.

\begin{theorem}\color{blue} \label{thmtype:9.2.2}
If $\omega \ne0$ and $m$ is a positive integer, then
\begin{align*}
\{e^{\lambda x}\cos\omega x, xe^{\lambda x}\cos\omega x,
&\dots, x^{m-1}e^{\lambda x}\cos\omega x,\\
e^{\lambda x}\sin\omega x, xe^{\lambda x}\sin\omega x,&
\dots, x^{m-1}e^{\lambda x}\sin\omega x\}
\end{align*}
is a fundamental set of solutions of
\[
[(D-\lambda)^2+\omega^2]^my=0.
\]
\end{theorem}

\begin{example}\label{example:9.2.6}
Find the general solution of
\begin{equation} \label{eq:9.2.19}
(D^2+4D+13)^3y=0.
\end{equation}
\end{example}

\begin{solution}
The characteristic polynomial of \eqref{eq:9.2.19} is
\[
p(r)=(r^2+4r+13)^3=\left((r+2)^2+9\right)^3.
\]
Therefore \eqref{eq:9.2.19} can be be written as
\[
[(D+2)^2+9]^3y=0,
\]
so \zcref{thmtype:9.2.2} implies that the general solution of
\eqref{eq:9.2.19} is
\[
y=(a_1+a_2x+a_3x^2)e^{-2x}\cos3x
+(b_1+b_2x+b_3x^2)e^{-2x}\sin3x.
\]
\end{solution}

\begin{example}\label{example:9.2.7}
Find the general solution of
\begin{equation} \label{eq:9.2.20}
y^{(4)}+4y'''+6y''+4y'=0.
\end{equation}
\end{example}

\begin{solution}
The characteristic polynomial of \eqref{eq:9.2.20} is
\begin{align*}
p(r)&=r^4+4r^3+6r^2+4r\\
&=r(r^3+4r^2+6r+4)\\
&=r(r+2)(r^2+2r+2)\\
&=r(r+2)[(r+1)^2+1].
\end{align*}
Therefore \eqref{eq:9.2.20} can be written as
\[
[(D+1)^2+1](D+2)Dy=0.
\]
Fundamental sets of solutions of
\[
\left[(D+1)^2+1\right] y=0,\quad (D+2) y=0,\quad\text{and}\quad  Dy=0.
\]
are given by
\[
\{e^{-x}\cos x,e^{-x}\sin x\},\quad \{e^{-2x}\},\quad\text{and}\quad  \{1\},
\]
respectively. Therefore the general solution of \eqref{eq:9.2.20} is
\[
y=e^{-x}(c_1\cos x+c_2\sin x)+c_3e^{-2x}+c_4.
\]
\end{solution}

\begin{example}\label{example:9.2.8}
Find a fundamental set of solutions of
\begin{equation} \label{eq:9.2.21}
[(D+1)^2+1]^2(D-1)^3(D+1)D^2y=0.
\end{equation}
\end{example}

\begin{solution}
A fundamental set of solutions of \eqref{eq:9.2.21} can be
obtained by combining fundamental sets of solutions of
\[
\left[(D+1)^2+1\right]^2 y=0,
\quad (D-1)^3 y=0,
\quad (D+1)y=0,
\quad\text{and}
\quad D^2y=0.
\]
Fundamental sets of solutions of these equations are given by
\[
\{e^{-x}\cos x, xe^{-x}\cos x, e^{-x}\sin x, xe^{-x}\sin x\},
\quad \{e^x, xe^x, x^2e^x\},%\\
\quad\{e^{-x}\},
\quad\text{and}
\quad\{1,x\},
\]
respectively. These ten functions form a fundamental set of solutions
of \eqref{eq:9.2.21}.
\end{solution}

 \exercises
 In \zcref[range]{exer:9.2.1,exer:9.2.14} find the general
solution.

\begin{exerciselistTwoColumn}

\item\label{exer:9.2.1}$y'''-3y''+3y'-y=0$
\item\label{exer:9.2.2}$y^{(4)}+8y''-9y=0$
\item\label{exer:9.2.3}$y'''-y''+16y'-16y=0$
\item\label{exer:9.2.4}$2y'''+3y''-2y'-3y=0$
\item\label{exer:9.2.5}$y'''+5y''+9y'+5y=0$%&
\item\label{exer:9.2.6}$4y'''-8y''+5y'-y=0$
\item\label{exer:9.2.7}$27y'''+27y''+9y'+y=0$%&
\item\label{exer:9.2.8}$y^{(4)}+y''=0$
\item\label{exer:9.2.9}$y^{(4)}-16y=0$%&
\item\label{exer:9.2.10}$y^{(4)}+12y''+36y=0$
\item\label{exer:9.2.11}$16y^{(4)}-72y''+81y=0$%&
\item\label{exer:9.2.12}$6y^{(4)}+5y'''+7y''+5y'+y=0$
\item(2)\label{exer:9.2.13}$4y^{(4)}+12y'''+3y''-13y'-6y=0$
\item(2)\label{exer:9.2.14}$y^{(4)}-4y'''+7y''-6y'+2y=0$
\end{exerciselistTwoColumn}
%
\textsl{In \zcref[range]{exer:9.2.15,exer:9.2.27} solve the
initial value problem. Where indicated by \CGex, graph the solution.}
\begin{exerciselist}
\item\label{exer:9.2.15}$y'''-2y''+4y'-8y=0, \quad  y(0)=2,\quad y'(0)=-2,\;  y''(0)=0$

\item\label{exer:9.2.16}$y'''+3y''-y'-3y=0, \quad  y(0)=0,\quad y'(0)=14,\quad y''(0)=-40$

\item\label{exer:9.2.17}\CGex
$y'''-y''-y'+y=0, \quad  y(0)=-2,\quad y'(0)=9,\quad y''(0)=4$

\item\label{exer:9.2.18}\CGex
$y'''-2y'-4y=0, \quad  y(0)=6,\quad y'(0)=3,\quad y''(0)=22$

\item\label{exer:9.2.19}\CGex
$3y'''-y''-7y'+5y=0, \quad  y(0)=\dfrac{14}{5},\quad y'(0)=0,\quad y''(0)=10$

\item\label{exer:9.2.20}$y'''-6y''+12y'-8y=0, \quad  y(0)=1,\quad y'(0)=-1,\quad y''(0)=-4$

\item\label{exer:9.2.21}$2y'''-11y''+12y'+9y=0, \quad  y(0)=6,\quad y'(0)=3,\quad y''(0)=13$

\item\label{exer:9.2.22}$8y'''-4y''-2y'+y=0, \quad  y(0)=4,\quad y'(0)=-3,\quad y''(0)=-1$

\item\label{exer:9.2.23}$y^{(4)}-16y=0, \quad   y(0)=2,\;  y'(0)=2,\;  y''(0)=-2,\;  y'''(0)=0$

\item\label{exer:9.2.24}$y^{(4)}-6y'''+7y''+6y'-8y=0, \quad  y(0)=-2,\quad y'(0)=-8,\quad y''(0)=-14$,

 $y'''(0)=-62$

\item\label{exer:9.2.25}$4y^{(4)}-13y''+9y=0, \quad  y(0)=1,\quad y'(0)=3,\quad y''(0)=1,\quad y'''(0)=3$

\item\label{exer:9.2.26}$y^{(4)}+2y'''-2y''-8y'-8y=0, \quad  y(0)=5,\quad y'(0)=-2,\quad y''(0)=6,\quad y'''(0)=8$

\item\label{exer:9.2.27}\CGex
$4y^{(4)}+8y'''+19y''+32y'+12y=0, \quad  y(0)=3,\quad y'(0)=-3,\quad y''(0)=
-\dfrac{7}{2}$,
 $y'''(0)=\dfrac{31}{4}$

\item\label{exer:9.2.28}Find a fundamental set of solutions of the given equation, and verify
that it's a fundamental set by evaluating  its Wronskian
at $x=0$.
\begin{alist*}[columns=2]
\item $(D-1)^2(D-2)y=0$
\item $(D^2+4)(D-3)y=0$
\item $(D^2+2D+2)(D-1)y=0$
\item $D^3(D-1)y=0$
\item $(D^2-1)(D^2+1)y=0$
\item $(D^2-2D+2)(D^2+1)y=0$
\end{alist*}
\end{exerciselist}
%
\textsl{In \zcref[range]{exer:9.2.29,exer:9.2.38} find a
fundamental set of solutions.}
\begin{exerciselistTwoColumn}
\item(2)\label{exer:9.2.29}$(D^2+6D+13)(D-2)^2D^3y=0$
\item(2)\label{exer:9.2.30}$(D-1)^2(2D-1)^3(D^2+1)y=0$
\item\label{exer:9.2.31}$(D^2+9)^3D^2y=0$%&
\item\label{exer:9.2.32}$(D-2)^3(D+1)^2Dy=0$
\item\label{exer:9.2.33}$(D^2+1)(D^2+9)^2(D-2)y=0$%&
\item\label{exer:9.2.34}$(D^4-16)^2y=0$
\item\label{exer:9.2.35}$(4D^2+4D+9)^3y=0$%&
\item\label{exer:9.2.36}$D^3(D-2)^2(D^2+4)^2y=0$
\item\label{exer:9.2.37}$(4D^2+1)^2(9D^2+4)^3y=0$%&
\item\label{exer:9.2.38}$\left[(D-1)^4-16\right]y=0$
\end{exerciselistTwoColumn}
\begin{exerciselist}
\item\label{exer:9.2.39}It can be shown that
\[
\begin{vmatrix}
1&1&\cdots&1\\
a_1&a_2&\cdots&a_n\\
a^2_1&a^2_2&\cdots&a^2_n\\
\vdots&\vdots&\ddots&\vdots\\
a^{n-1}_1&a^{n-1}_2&\cdots&a^{n-1}_n
\end{vmatrix}=
\prod_{1\le i<j\le n}(a_j-a_i),
\tag{A}
\]
where the left side is  the
\personHref{Vandermonde}{Vandermonde determinant}\index{Vandermonde determinant} and
the right side is the product of all factors of the form $(a_j-a_i)$
with $i$ and $j$ between $1$ and $n$ and $i<j$.

\begin{alist}
\item %(a)
Verify  (A) for $n=2$ and $n=3$.

\item %(b)
Find the Wronskian of $\{e^{{a_1}x}, \quad  e^{{a_2}x},\dots, e^{{a_n}x}\}$.
\end{alist}

\item\label{exer:9.2.40}A theorem from algebra says that if $P_1$ and $P_2$ are polynomials
with no common factors then there are polynomials $Q_1$ and $Q_2$
such that
\[
Q_1P_1+Q_2P_2=1.
\]
This implies that
\[
Q_1(D)P_1(D)y+Q_2(D)P_2(D)y=y
\]
for every function $y$ with enough derivatives for the left side to be
defined.

\begin{alist}
\item %(a)
Use this to show that if $P_1$ and $P_2$ have no common factors and
\[
P_1(D)y=P_2(D)y=0
\]
then $y=0$.

\item %(b)
Suppose $P_1$ and $P_2$ are polynomials with no common factors.
Let $u_1$, \dots, $u_r$ be linearly independent solutions of $P_1(D)y=0$
and let $v_1$, \dots, $v_s$ be linearly independent solutions of
$P_2(D)y=0$. Use (a) to show that $\{u_1,\dots,u_r,\allowbreak
v_1,\dots,v_s\}$ is a linearly independent set.

\item %(c)
Suppose the characteristic polynomial of the constant coefficient
equation
\[
a_0y^{(n)}+a_1y^{(n-1)}+\cdots+a_ny=0
\tag{A}
\]
has the factorization
\[
p(r)=a_0p_1(r)p_2(r)\cdots p_k(r),
\]
where each $p_j$ is of the form
\[
p_j(r)=(r-r_j)^{n_j} \quad\text{or}\quad
p_j(r)=[(r-\lambda_j)^2+w^2_j]^{m_j}\qquad  (\omega_j>0)
\]
and no two of the polynomials $p_1$, $p_2$, \dots, $p_k$ have a common
factor. Show that we can find a fundamental set of solutions
$\{y_1,y_2,\dots,y_n\}$ of (A) by finding a
fundamental set of solutions of each of the equations
\[
p_j(D)y=0,\quad 1\le j\le k,
\]
 and taking $\{y_1,y_2,\dots,y_n\}$ to be the set of all
functions in these separate fundamental sets.
\end{alist}

\item\label{exer:9.2.41}\begin{alist}
\item % (a)
Show that if
\[
z=p(x)\cos\omega x+q(x)\sin\omega x,
\tag{A}
\]
where $p$ and $q$ are polynomials of degree $\le k$, then
\[
(D^2+\omega^2)z=p_1(x)\cos\omega x+q_1(x)\sin\omega x,
\]
where $p_1$ and $q_1$ are polynomials of degree $\le k-1$.

\item % (b)
Apply (a) $m$ times to show that if $z$ is of
the form (A) where $p$ and $q$ are polynomial of
degree $\le m-1$, then
\[
(D^2+\omega^2)^mz=0.
\tag{B}
\]

\item %(c)
Use Eqn.~\eqref{eq:9.2.17} to show that if $y=e^{\lambda x}z$ then
\[
[(D-\lambda)^2+\omega^2]^my=e^{\lambda
x}(D^2+\omega^2)^mz.
\]

\item % (d)
Conclude from (b) and (c) that if $p$ and $q$ are arbitrary
polynomials of degree $\le m-1$ then
\[
y=e^{\lambda x}(p(x)\cos\omega x+q(x)\sin\omega x)
\]
is a solution of
\[
[(D-\lambda)^2+\omega^2]^my=0.
\tag{C}
\]

\item % (e)
Conclude from (d) that the functions
\[
\begin{aligned}
e^{\lambda x}\cos\omega x, xe^{\lambda x}\cos\omega x,
&\dots, x^{m-1}e^{\lambda x}\cos\omega x,\\
e^{\lambda x}\sin\omega x, xe^{\lambda x}\sin\omega x,&
\dots, x^{m-1}e^{\lambda x}\sin\omega x
\end{aligned}
\tag{D}
\]
are all solutions of (C).

\item % (f)
Complete the proof of \zcref{thmtype:9.2.2} by showing that the
functions in (D) are linearly independent.
\end{alist}

\item\label{exer:9.2.42}\begin{alist}
\item % (a)
Use the trigonometric identities
\begin{align*}
\cos(A+B)&=\cos A\cos B-\sin A\sin B\\
\sin(A+B)&=\cos A\sin B+\sin A\cos B
\end{align*}
to show that
\[
(\cos A+i\sin A)(\cos B+i\sin B)=\cos(A+B)+i\sin(A+B).
\]
\item % (b)
Apply (a) repeatedly to show that if $n$ is a positive
integer then
\[
\prod_{k=1}^n(\cos A_k+i\sin A_k)=\cos(A_1+A_2+\cdots+A_n)
+i\sin(A_1+A_2+\cdots+A_n).
\]
\item % (c)
Infer from (b) that if $n$ is a positive integer then
\[
(\cos\theta+i\sin\theta)^n=\cos n\theta+i\sin n\theta.
\tag{A}
\]
\item % (d)
Show that (A) also holds if $n=0$ or a negative integer.
\hint{Verify by direct calculation that
\[
(\cos\theta+i\sin\theta)^{-1}=(\cos\theta-i\sin\theta).
\]
Then replace $\theta$ by $-\theta$  in (A).}
\item % (e)
Now suppose $n$ is a positive integer. Infer from (A)
that if
\[
z_k=\cos\left(\frac{2k\pi}{n}\right)+i\sin\left(\frac{2k\pi}{n}\right)
,\quad k=0,1,\dots,n-1,
\]
and
\[
\zeta_k=\cos\left(\frac{(2k+1)\pi}{n}\right)+i\sin\left(\frac{(2k+1)\pi}{n}\right)
,\quad k=0,1,\dots,n-1,
\]
then
\[
z_k^n=1\quad\text{and}\quad\zeta_k^n=-1,\quad k=0,1,\dots,n-1.
\]
(Why don't we also consider other integer values for $k$?)

\item % (f)
Let $\rho$ be a positive number. Use (e)  to show that
\[
z^n-\rho=(z-\rho^{1/n} z_0)(z-\rho^{1/n}z_1)\cdots(z-\rho^{1/n} z_{n-1})
\]
and
\[
z^n+\rho=(z-\rho^{1/n} \zeta_0)(z-\rho^{1/n} \zeta_1)\cdots(z-\rho^{1/n}
\zeta_{n-1}).
\]
\end{alist}

\item\label{exer:9.2.43}Use (e) of \zcref{exer:9.2.42} to find a fundamental set of
solutions of the given equation.
\begin{alist*}[columns=2]
\item $y'''-y=0$
\item $y'''+y=0$
\item $y^{(4)}+64y=0$
\item $y^{(6)}-y=0$
\item $y^{(6)}+64y=0$
\item $\left[(D-1)^6-1\right]y=0$
\item $y^{(5)}+y^{(4)}+y'''+y''+y'+y=0$
\end{alist*}

\item\label{exer:9.2.44}An equation of the form
\[
 a_0x^ny^{(n)}+a_1x^{n-1}y^{(n-1)}+\cdots
+a_{n-1}xy'+a_ny=0,\quad x>0,
\tag{A}
\]
where $a_0$, $a_1$, \dots, $a_n$ are constants, is an \emph{\color{blue}Euler} or
\emph{\color{blue}equidimensional} equation\index{equidimensional equation}.

Show that if
\[
 x=e^t \quad \text{and} \quad Y(t)=y(x(t)),
\tag{B}
\]
then
\begin{align*}
x \frac{dy}{dx}&=\frac{dY}{dt}\\
x^2\frac{d^2y}{dx^2}&=\frac{d^2Y}{dt^2}-\frac{dY}{dt}\\
x^3\frac{d^3y}{dx^3}&=\frac{d^3Y}{dt^3}-3\frac{d^2Y}{dt^2}+2\frac{dY}{dt}.
\end{align*}
In general, it can be shown that if $r$ is any integer $\ge2$ then
\[
x^r \frac{d^ry}{dx^r}=\frac{d^rY}{dt^r}+
 A_{1r}\frac{d^{r-1}Y}{dt^{r-1}}+\cdots+A_{r-1,r}
 \frac{dY}{dt}
\]
where $A_{1r}$, \dots, $A_{r-1,r}$ are integers.  Use these results to
show that the substitution (B) transforms (A) into a
constant coefficient equation for $Y$ as a function of $t$.


\item\label{exer:9.2.45}Use \zcref{exer:9.2.44}  to show that a function $y=y(x)$ satisfies
the equation
\[
 a_0x^3y'''+a_1x^2y''+a_2xy'+a_3y=0,
\tag{A}
\]
on $(0,\infty)$
if and only if the function $Y(t)=y(e^t)$ satisfies
\[
a_0\frac{d^3Y}{dt^3}+(a_1-3a_0) \frac{d^2Y}{dt^2}+(a_2-a_1+2a_0)
\frac{dY}{dt}+a_3Y=0.
\]
Assuming that $a_0$, $a_1$, $a_2$, $a_3$ are real and $a_0 \ne0$, find
the possible forms for the general solution of (A).
\end{exerciselist}



\section{Undetermined Coefficients for Higher Order  Equations}
\label{sec:higher_undetermined}

In this section we consider the constant coefficient equation
\begin{equation}
 \label{eq:9.3.1} a_0y^{(n)}+a_1y^{(n-1)}+\cdots+a_ny=F(x),
\end{equation}
where $n\ge3$ and $F$ is a linear combination of
functions of the form
\[
e^{\alpha x}\left(p_0+p_1x+\cdots+p_kx^k\right)
\]
or
\[
e^{\lambda x}\left[\left(p_0+p_1x+\cdots+p_kx^k\right)
\cos\omega x+
\left(q_0+q_1x+\cdots+q_kx^k\right)
\sin\omega x\right].
\]

From \zcref{thmtype:9.1.5}, the general solution of \eqref{eq:9.3.1} is
$y=y_p+y_c$, where $y_p$ is a particular solution of \eqref{eq:9.3.1} and
$y_c$ is the general solution of the complementary equation
\[
a_0y^{(n)}+a_1y^{(n-1)}+\cdots+a_ny=0.
\]
In \zcref{sec:higher_const_coeff} we learned how to find $y_c$. Here
we
will learn how to find $y_p$ when the forcing function has the form stated
above. The procedure that we  use is a generalization of the
method that we used in
\zcref{sec:undetermined_coeffs_i,sec:undetermined_coeffs_ii},
and is again called
\emph{\color{blue}method of undetermined coefficients}. Since
the underlying ideas are the same as those in
\zcref*{sec:undetermined_coeffs_i,sec:undetermined_coeffs_ii},
we'll give  an informal presentation based on
examples.


\boxit{Forcing Functions of the Form
$e^{\alpha x}\left(p_0+p_1x+\cdots+p_kx^k\right)$}

We  first consider  equations of the form
\[
a_0y^{(n)}+a_1y^{(n-1)}+\cdots+a_ny=e^{\alpha
x}\left(p_0+p_1x+\cdots+p_kx^k\right).
\]

\begin{example}\label{example:9.3.1}
Find a particular solution of
\begin{equation} \label{eq:9.3.2}
y'''+3y''+2y'-y=e^x(21+24x+28x^2+5x^3).
\end{equation}
\end{example}

\begin{solution}
Substituting
\begin{align*}
y&=ue^x,\\
y'&=e^x(u'+u),\\
y''&=e^x(u''+2u'+u),\\
y'''&=e^x(u'''+3u''+3u'+u)
\end{align*}
into \eqref{eq:9.3.2} and canceling $e^x$ yields
\[
(u'''+3u''+3u'+u)+3(u''+2u'+u)+2(u'+u)-u
=21+24x+28x^2+5x^3,
\]
or
\begin{equation} \label{eq:9.3.3}
u'''+6u''+11u'+5u=21+24x+28x^2+5x^3.
\end{equation}
Since the unknown $u$ appears on the left, we can see that
\eqref{eq:9.3.3} has a particular solution of the form
\[
u_p=A+Bx+Cx^2+Dx^3.
\]
Then
\begin{align*}
u_p'&=B+2Cx+3Dx^2\\
u_p''&=2C+6Dx\\
u_p'''&=6D.
\end{align*}
Substituting from the last four equations  into the left side of
\eqref{eq:9.3.3} yields
\begin{align*}
u_p'''+6u_p''+11u_p'+5u_p
&=6D+6(2C+6Dx)+11(B+2Cx+3Dx^2)\\
&\qquad{}+5(A+Bx+Cx^2+Dx^3)\\
&=(5A+11B+12C+6D)+(5B+22C+36D)x\\
&\qquad{}+(5C+33D)x^2+5Dx^3.
\end{align*}
Comparing coefficients of like powers of $x$ on the right sides of
this equation and \eqref{eq:9.3.3} shows that $u_p$ satisfies \eqref{eq:9.3.3}
if
\begin{align*}
5D&=5\\
5C+33D&=28\\
5B+22C+36D&=24\\
5A+11B+12C+\phantom{3}6D&=21.
\end{align*}
Solving these equations successively yields $D=1$, $C=-1$, $B=2$, $A=1$.
Therefore
\[
u_p=1+2x-x^2+x^3
\]
is a particular solution of  \eqref{eq:9.3.3}, so
\[
y_p=e^xu_p=e^x(1+2x-x^2+x^3)
\]
is a particular solution of  \eqref{eq:9.3.2} (\zcref{figure:9.3.1}).
\end{solution}

\begin{figure}
  \centering\color{blue}
  \includegraphics[alt={A plot of the solution for 0<x<2 and 0<y<50 that begins at (0,1), increases slowly at first, and then increases rapidly.},
  width=5.67in,height=3.3in,keepaspectratio]{fig090301}
  \caption{$y_p=e^x(1+2x-x^2+x^3)$ }
  \label{figure:9.3.1}
\end{figure}


\begin{example}\label{example:9.3.2}
Find a particular solution of
\begin{equation} \label{eq:9.3.4}
y^{(4)}-y'''-6y''+4y'+8y=e^{2x}(4+19x+6x^2).
\end{equation}
\end{example}

\begin{solution}
Substituting
\begin{align*}
y&=ue^{2x},\\
y'&=e^{2x}(u'+2u),\\
y''&=e^{2x}(u''+4u'+4u),\\
y'''&=e^{2x}(u'''+6u''+12u'+8u),\\
y^{(4)}&=e^{2x}(u^{(4)}+8u'''+24u''+32u'+16u)
\end{align*}
into \eqref{eq:9.3.4} and canceling $e^{2x}$ yields
\begin{multline*}
(u^{(4)}+8u'''+24u''+32u'+16u)-(u'''+6u''+12u'+8u)\\
-6(u''+4u'+4u)+4(u'+2u)+8u=4+19x+6x^2,
\end{multline*}
or
\begin{equation} \label{eq:9.3.5}
u^{(4)}+7u'''+12u''=4+19x+6x^2.
\end{equation}
Since neither $u$ nor $u'$ appear on the left, we can see that
\eqref{eq:9.3.5} has a particular solution of the form
\begin{equation} \label{eq:9.3.6}
u_p=Ax^2+Bx^3+Cx^4.
\end{equation}
Then
\begin{align*}
u_p'&=2Ax+3Bx^2+4Cx^3\\
u_p''&=2A+6Bx+12Cx^2\\
u_p'''&=6B+24Cx\\
u_p^{(4)}&=24C.
\end{align*}
Substituting $u_p''$, $u_p'''$, and $u_p^{(4)}$  into the left side of
\eqref{eq:9.3.5} yields
\begin{align*}
u_p^{(4)}+7u_p'''+12u_p''
&=24C+7(6B+24Cx)+12(2A+6Bx+12Cx^2)\\
&=(24A+42B+24C)+(72B+168C)x+144Cx^2.
\end{align*}
Comparing coefficients of like powers of $x$ on the right sides of
this equation and \eqref{eq:9.3.5} shows that $u_p$ satisfies \eqref{eq:9.3.5}
if
\begin{align*}
144C&=6\\
72B+168C&=19\\
24A+42B+\phantom{1}24C&=4.
\end{align*}
Solving these equations successively yields $C=1/24$, $B=1/6$, $A=-1/6$.
Substituting these into \eqref{eq:9.3.6} shows that
\[
u_p=\frac{x^2}{24}(-4+4x+x^2)
\]
is a particular solution of  \eqref{eq:9.3.5}, so
\[
y_p=e^{2x}u_p=\frac{x^2e^{2x}}{24}(-4+4x+x^2)
\]
is a particular solution of  \eqref{eq:9.3.4}.
(\zcref{figure:9.3.2}).
\end{solution}

\begin{figure}
  \centering\color{blue}
  \includegraphics[alt={A graph of the solution curve for 0<x<1 and -0.1<y<0.3, beginning flat at the origin, bottoming out near (0.65,-0.05), crossing the x axis near x=0.83, and increasing from there.},
  width=5.67in,height=3.3in,keepaspectratio]{fig090302}
  \caption{$y_p=\dfrac{x^2e^{2x}}{24}(-4+4x+x^2)$}
  \label{figure:9.3.2}
\end{figure}


\boxit{Forcing Functions of the Form
$e^{\lambda x}\left(P(x)\cos\omega x+Q(x)\sin\omega x\right)$}

We  now consider equations of the form
\[
a_0y^{(n)}+a_1y^{(n-1)}+\cdots+a_ny=
e^{\lambda x}\left(P(x)\cos\omega x+Q(x)\sin\omega x\right),
\]
where $P$ and $Q$ are polynomials.

\begin{example}\label{example:9.3.3}
Find a particular solution of
\begin{equation} \label{eq:9.3.7}
y'''+y''-4y'-4y=e^x[(5-5x)\cos x+(2+5x)\sin x].
\end{equation}
\end{example}

\begin{solution}
Substituting
\begin{align*}
y&=ue^x,\\
y'&=e^x(u'+u),\\
y''&=e^x(u''+2u'+u),\\
y'''&=e^x(u'''+3u''+3u'+u)
\end{align*}
into \eqref{eq:9.3.7} and canceling $e^x$ yields
\[
(u'''+3u''+3u'+u)+(u''+2u'+u)-4(u'+u)-4u
=(5-5x)\cos x+(2+5x)\sin x,
\]
or
\begin{equation} \label{eq:9.3.8}
u'''+4u''+u'-6u=(5-5x)\cos x+(2+5x)\sin x.
\end{equation}
Since $\cos x$ and $\sin x$ are not solutions of the complementary
equation
\[
u'''+4u''+u'-6u=0,
\]
a theorem analogous to \zcref{thmtype:5.5.1} implies that
\eqref{eq:9.3.8} has a particular solution of the form
\begin{equation} \label{eq:9.3.9}
u_p=(A_0+A_1x)\cos x+(B_0+B_1x)\sin x.
\end{equation}
Then
\begin{align*}
u_p'&=(A_1+B_0+B_1x)\cos x+(B_1-A_0-A_1x)\sin x,\\
u_p''&=(2B_1-A_0-A_1x)\cos x-(2A_1+B_0+B_1x)\sin x,\\
u_p'''&=-(3A_1+B_0+B_1x)\cos x-(3B_1-A_0-A_1x)\sin x,\\
\end{align*}
so
\begin{align*}
u_p'''+4u_p''+u_p'-6u_p
&=-\left[10A_0+2A_1-8B_1+10A_1x\right]\cos x\\
&\qquad{}- \left[10B_0+2B_1+8A_1+10B_1x\right]\sin x.
\end{align*}
Comparing the coefficients of $x\cos x$, $x\sin x$, $\cos x$, and
$\sin x$ here with the corresponding coefficients in \eqref{eq:9.3.8}
shows that $u_p$ is a solution of \eqref{eq:9.3.8} if
\begin{align*}
-10A_1&=-5\\
-10B_1&=\phantom{-}5\\
-10A_0-2A_1+8B_1&=\phantom{-}5\\
-10B_0-2B_1-8A_1&=\phantom{-}2.
\end{align*}
Solving the first two equations yields $A_1=1/2$, $B_1=-1/2$.
Substituting these into the last two equations  yields
\begin{align*}
-10A_0&=5+2A_1-8B_1=10\\
-10B_0&=2+2B_1+8A_1=5,
\end{align*}
so $A_0=-1$, $B_0=-1/2$.
Substituting $A_0=-1$, $A_1=1/2$, $B_0=-1/2$, $B_1=-1/2$ into
\eqref{eq:9.3.9} shows that
\[
u_p=-\frac{1}{2}\left[(2-x)\cos
x+(1+x)\sin x\right]
\]
is a particular solution of  \eqref{eq:9.3.8}, so
\[
y_p=e^xu_p=-\frac{e^x}{2}\left[(2-x)\cos x+(1+x)\sin x\right]
\]
is a particular  solution of \eqref{eq:9.3.7} (\zcref{figure:9.3.3}).
\end{solution}

\begin{figure}
  \centering\color{blue}
  \includegraphics[alt={A graph of the solution curve for 0<x<4 and -20<y<100, beginning flat at (0,-1), bottoming out near (2.8,-15), crossing the x axis near x=3.4, and increasing from there.},
  width=5.67in,height=3.3in,keepaspectratio]{fig090303}
\caption{$y_p=e^xu_p=-\dfrac{e^x}{2}\left[(2-x)\cos x+(1+x)\sin x\right]$}
  \label{figure:9.3.3}
\end{figure}


\begin{example}\label{example:9.3.4}
Find a particular solution of
\begin{equation} \label{eq:9.3.10}
 y'''+4y''+6y'+4y=
e^{-x}\left[(1-6x)\cos x-(3+2x)\sin x\right].
\end{equation}
\end{example}

\begin{solution}
Substituting
\begin{align*}
y&=ue^{-x},\\
y'&=e^{-x}(u'-u),\\
y''&=e^{-x}(u''-2u'+u),\\
y'''&=e^{-x}(u'''-3u''+3u'-u)
\end{align*}
into \eqref{eq:9.3.10} and canceling $e^{-x}$ yields
\[
(u'''-3u''+3u'-u)+4(u''-2u'+u)+6(u'-u)+4u
=(1-6x)\cos x-(3+2x)\sin x,
\]
or
\begin{equation} \label{eq:9.3.11}
u'''+u''+u'+u=(1-6x)\cos x-(3+2x)\sin x.
\end{equation}
Since $\cos x$ and $\sin x$ are  solutions of the complementary
equation
\[
u'''+u''+u'+u=0,
\]
a theorem analogous to \zcref{thmtype:5.5.1} implies
that  \eqref{eq:9.3.11} has a particular solution of the form
\begin{equation} \label{eq:9.3.12}
u_p=(A_0x+A_1x^2)\cos x+(B_0x+B_1x^2)\sin x.
\end{equation}

Then
\begin{align*}
u_p'&=[A_0+(2A_1+B_0)x+B_1x^2]\cos x+[B_0+(2B_1-A_0)x-A_1x^2]\sin x,\\
u_p''&=[2A_1+2B_0-(A_0-4B_1)x-A_1x^2]\cos x
+[2B_1-2A_0-(B_0+4A_1)x-B_1x^2]\sin x,\\
u_p'''&=-[3A_0-6B_1+(6A_1+B_0)x+B_1x^2]\cos x
-[3B_0+6A_1+(6B_1-A_0)x-A_1x^2]\sin x,
\end{align*}
so
\begin{align*}
u_p'''+u_p''+u_p'+u_p
&=-[2A_0-2B_0-2A_1-6B_1+(4A_1-4B_1)x]\cos x\\
&\qquad{}-[2B_0+2A_0-2B_1+6A_1+(4B_1+4A_1)x]\sin x.
\end{align*}
Comparing the coefficients of $x\cos x$, $x\sin x$, $\cos x$, and
$\sin x$ here with the corresponding coefficients in \eqref{eq:9.3.11}
shows that $u_p$ is a solution of \eqref{eq:9.3.11} if
\begin{align*}
-4A_1+4B_1&=-6\\
-4A_1-4B_1&=-2\\
-2A_0+2B_0+2A_1+6B_1&=\phantom{-}1\\
-2A_0-2B_0-6A_1+2B_1&=-3.
\end{align*}
Solving the first two equations yields $A_1=1$, $B_1=-1/2$.
Substituting these into the last two equations  yields
\begin{align*}
-2A_0+2B_0&=\phantom{-}1-2A_1-6B_1=2\\
-2A_0-2B_0&=-3+6A_1-2B_1=4,
\end{align*}
%
\begin{figure}
  \centering\color{blue}
  \includegraphics[alt={A wobbly graph of the solution, beginning at the origin, decreasing to a minimum near (2,-0.42), crossing the x axis near x=3.9, peaking near (5,0.12), crossing the x axis again near x=7.2, and then increasing to the horizontal asymptote y=0.},
  width=5.67in,height=3.3in,keepaspectratio]{fig090304}
\caption{ $y_p=-\dfrac{xe^{-x}}{2}\left[(3-2x)\cos x+(1+x)\sin x\right]$}
  \label{figure:9.3.4}
\end{figure}
%
so $A_0=-3/2$ and $B_0=-1/2$.
Substituting $A_0=-3/2$, $A_1=1$, $B_0=-1/2$, $B_1=-1/2$ into
\eqref{eq:9.3.12} shows that
\[
u_p=-\frac{x}{2}\left[(3-2x)\cos
x+(1+x)\sin x\right]
\]
is a particular solution of  \eqref{eq:9.3.11}, so
\[
 y_p=e^{-x}u_p=-\frac{xe^{-x}}{2}\left[(3-2x)\cos
x+(1+x)\sin x\right]
\]
(\zcref{figure:9.3.4}) is a particular  solution of \eqref{eq:9.3.10}.
\end{solution}



\exercises
In \zcref[range]{exer:9.3.1,exer:9.3.59} find a particular solution.

\begin{exerciselist}

\item\label{exer:9.3.1}$y'''-6y''+11y'-6y=-e^{-x}(4+76x-24x^2)$
\item\label{exer:9.3.2}$y'''-2y''-5y'+6y=e^{-3x}(32-23x+6x^2)$
\item\label{exer:9.3.3}$4y'''+8y''-y'-2y=-e^x(4+45x+9x^2)$
\item\label{exer:9.3.4}$y'''+3y''-y'-3y=e^{-2x}(2-17x+3x^2)$
\item\label{exer:9.3.5}$y'''+3y''-y'-3y=e^x(-1+2x+24x^2+16x^3)$
\item\label{exer:9.3.6}$y'''+y''-2y=e^x(14+34x+15x^2)$
\item\label{exer:9.3.7}$4y'''+8y''-y'-2y=-e^{-2x}(1-15x)$
\item\label{exer:9.3.8}$y'''-y''-y'+y=e^x(7+6x)$
\item\label{exer:9.3.9}$2y'''-7y''+4y'+4y=e^{2x}(17+30x)$
\item\label{exer:9.3.10}$y'''-5y''+3y'+9y=2e^{3x}(11-24x^2)$
\item\label{exer:9.3.11}$y'''-7y''+8y'+16y=2e^{4x}(13+15x)$
\item\label{exer:9.3.12}$8y'''-12y''+6y'-y=e^{x/2}(1+4x)$
\item\label{exer:9.3.13}$y^{(4)}+3y'''-3y''-7y'+6y=-e^{-x}(12+8x-8x^2)$
\item\label{exer:9.3.14}$y^{(4)}+3y'''+y''-3y'-2y=-3e^{2x}(11+12x)$
\item\label{exer:9.3.15}$y^{(4)}+8y'''+24y''+32y'=-16e^{-2x}(1+x+x^2-x^3)$
\item\label{exer:9.3.16}$4y^{(4)}-11y''-9y'-2y=-e^x(1-6x)$
\item\label{exer:9.3.17}$y^{(4)}-2y'''+3y'-y=e^x(3+4x+x^2)$
\item\label{exer:9.3.18}$y^{(4)}-4y'''+6y''-4y'+2y=e^{2x}(24+x+x^4)$
\item\label{exer:9.3.19}$2y^{(4)}+5y'''-5y'-2y=18e^x(5+2x)$
\item\label{exer:9.3.20}$y^{(4)}+y'''-2y''-6y'-4y=-e^{2x}(4+28x+15x^2)$
\item\label{exer:9.3.21}$2y^{(4)}+y'''-2y'-y=3e^{-x/2}(1-6x)$
\item\label{exer:9.3.22}$y^{(4)}-5y''+4y=e^x(3+x-3x^2)$
\item\label{exer:9.3.23}$y^{(4)}-2y'''-3y''+4y'+4y=e^{2x}(13+33x+18x^2)$
\item\label{exer:9.3.24}$y^{(4)}-3y'''+4y'=e^{2x}(15+26x+12x^2)$
\item\label{exer:9.3.25}$y^{(4)}-2y'''+2y'-y=e^x(1+x)$
\item\label{exer:9.3.26}$2y^{(4)}-5y'''+3y''+y'-y=e^x(11+12x)$
\item\label{exer:9.3.27}$y^{(4)}+3y'''+3y''+y'=e^{-x}(5-24x+10x^2)$
\item\label{exer:9.3.28}$y^{(4)}-7y'''+18y''-20y'+8y=e^{2x}(3-8x-5x^2)$
\item\label{exer:9.3.29}$y'''-y''-4y'+4y=e^{-x}\left[(16+10x)\cos
x+(30-10x)\sin x\right]$
\item\label{exer:9.3.30}$y'''+y''-4y'-4y=e^{-x}\left[(1-22x)\cos
2x-(1+6x)\sin2x\right]$
\item\label{exer:9.3.31}$y'''-y''+2y'-2y=e^{2x}[(27+5x-x^2)\cos
x+(2+13x+9x^2)\sin x]$
\item\label{exer:9.3.32}$y'''-2y''+y'-2y=-e^x[(9-5x+4x^2)\cos
2x-(6-5x-3x^2)\sin2x]$
\item\label{exer:9.3.33}$y'''+3y''+4y'+12y=8\cos2x-16\sin2x$
\item\label{exer:9.3.34}$y'''-y''+2y=e^x[(20+4x)\cos x-(12+12x)\sin
 x]$
\item\label{exer:9.3.35}$y'''-7y''+20y'-24y=-e^{2x}[(13-8x)\cos
2x-(8-4x)\sin2x]$
\item\label{exer:9.3.36}$y'''-6y''+18y'=-e^{3x}[(2-3x)\cos
3x-(3+3x)\sin3x]$
\item\label{exer:9.3.37}$y^{(4)}+2y'''-2y''-8y'-8y=e^x(8\cos
x+16\sin x)$
\item\label{exer:9.3.38}$y^{(4)}-3y'''+2y''+2y'-4y=e^x(2\cos2x
-\sin2x)$
\item\label{exer:9.3.39}$y^{(4)}-8y'''+24y''-32y'+15y=e^{2x}(15x\cos2x+32\sin2x)$
\item\label{exer:9.3.40}$y^{(4)}+6y'''+13y''+12y'+4y=e^{-x}[(4-x)\cos x-(5+x)\sin x]$
\item\label{exer:9.3.41}$y^{(4)}+3y'''+2y''-2y'-4y=-e^{-x}
(\cos x-\sin x)$
\item\label{exer:9.3.42}$y^{(4)}-5y'''+13y''-19y'+10y=e^x
(\cos2x+\sin2x)$
\item\label{exer:9.3.43}$y^{(4)}+8y'''+32y''+64y'+39y=e^{-2x}[(4-15x)\cos3x-(4+15x)\sin
3x]$
\item\label{exer:9.3.44}$y^{(4)}-5y'''+13y''-19y'+10y=e^x[(7+8x)\cos
2x+(8-4x)\sin2x]$
\item\label{exer:9.3.45}$y^{(4)}+4y'''+8y''+8y'+4y=-2e^{-x}
(\cos x-2\sin x)$
\item\label{exer:9.3.46}$y^{(4)}-8y'''+32y''-64y'+64y=e^{2x}
(\cos2x-\sin2x)$
\item\label{exer:9.3.47}$y^{(4)}-8y'''+26y''-40y'+25y=e^{2x}[3\cos
x-(1+3x)\sin x]$
\item\label{exer:9.3.48}$y'''-4y''+5y'-2y=e^{2x}-4e^x-2\cos x+4\sin x$
\item\label{exer:9.3.49}$y'''-y''+y'-y=5e^{2x}+2e^x-4\cos x+4\sin x$
\item\label{exer:9.3.50}$y'''-y'=-2(1+x)+4e^x-6e^{-x}+96e^{3x}$
\item\label{exer:9.3.51}$y'''-4y''+9y'-10y=10e^{2x}+20e^x\sin2x-10$
\item\label{exer:9.3.52}$y'''+3y''+3y'+y=12e^{-x}+9\cos2x-13\sin2x$
\item\label{exer:9.3.53}$y'''+y''-y'-y=4e^{-x}(1-6x)-2x\cos x+2(1+x)\sin
x $
\item\label{exer:9.3.54}$y^{(4)}-5y''+4y=-12e^x+6e^{-x}+10\cos x$
\item\label{exer:9.3.55}$y^{(4)}-4y'''+11y''-14y'+10y=-e^x(\sin
x+2\cos2x)$
\item\label{exer:9.3.56}$y^{(4)}+2y'''-3y''-4y'+4y=2e^x(1+x)+e^{-2x}$
\item\label{exer:9.3.57}$y^{(4)}+4y=\sinh x\cos x-\cosh x\sin x$
\item\label{exer:9.3.58}$y^{(4)}+5y'''+9y''+7y'+2y=e^{-x}(30+24x)-e^{-2x}$
\item\label{exer:9.3.59}$y^{(4)}-4y'''+7y''-6y'+2y=e^x(12x-2\cos
x+2\sin x)$

\exercisetext{
In \zcref[range]{exer:9.3.60,exer:9.3.68} find the general solution.}

\item\label{exer:9.3.60}$y'''-y''-y'+y=e^{2x}(10+3x)$
\item\label{exer:9.3.61}$y'''+y''-2y=-e^{3x}(9+67x+17x^2)$
\item\label{exer:9.3.62}$y'''-6y''+11y'-6y=e^{2x}(5-4x-3x^2)$
\item\label{exer:9.3.63}$y'''+2y''+y'=-2e^{-x}(7-18x+6x^2)$
\item\label{exer:9.3.64}$y'''-3y''+3y'-y=e^x(1+x)$
\item\label{exer:9.3.65}$y^{(4)}-2y''+y=-e^{-x}(4-9x+3x^2)$
\item\label{exer:9.3.66}$y'''+2y''-y'-2y=e^{-2x}\left[(23-2x)\cos
x+(8-9x)\sin x\right]$
\item\label{exer:9.3.67}$y^{(4)}-3y'''+4y''-2y'=e^x\left[(28+6x)\cos
2x+(11-12x)\sin2x\right]$
\item\label{exer:9.3.68}$y^{(4)}-4y'''+14y''-20y'+25y=e^x\left[(2+6x)\cos
2x+3\sin2x\right]$


\exercisetext{
In \zcref[range]{exer:9.3.69,exer:9.3.74} solve the initial value
problem and graph the solution.}


\item\label{exer:9.3.69}\CGex\  $y'''-2y''-5y'+6y=2e^x(1-6x),\quad
y(0)=2,
\quad y'(0)=7,\quad y''(0)=9$
\item\label{exer:9.3.70}\CGex\  $y'''-y''-y'+y=-e^{-x}(4-8x),\quad
y(0)=2,
\quad y'(0)=0,\quad y''(0)=0$
\item\label{exer:9.3.71}\CGex\  $4y'''-3y'-y=e^{-x/2}(2-3x),\quad
y(0)=-1,
\quad y'(0)=15,\quad y''(0)=-17$
\item\label{exer:9.3.72}\CGex\
$y^{(4)}+2y'''+2y''+2y'+y=e^{-x}(20-12x),\,
y(0)=3,\; y'(0)=-4,\; y''(0)=7,\; y'''(0)=-22$
\item\label{exer:9.3.73}\CGex\ $y'''+2y''+y'+2y=30\cos x-10\sin x,
\quad y(0)=3,\quad y'(0)=-4,\quad y''(0)=16$
\item\label{exer:9.3.74}\CGex\ $y^{(4)}-3y'''+5y''-2y'=-2e^x(\cos x-\sin
x),\; y(0)=2,\; y'(0)=0,\; y''(0)~=~-1, \\ \hspace*{30pt} y'''(0)=-5$

\item\label{exer:9.3.75}Prove:
A function $y$ is a solution of the constant coefficient
nonhomogeneous equation
\[
a_0y^{(n)}+a_1y^{(n-1)}+\cdots+a_ny=e^{\alpha x}G(x)
\tag{A}
\]
if and only if $y=ue^{\alpha x}$, where $u$ satisfies the differential
equation
\[
a_0u^{(n)}+\frac{p^{(n-1)}(\alpha)}{(n-1)!}u^{(n-1)}+
\frac{p^{(n-2)}(\alpha)}{(n-2)!}u^{(n-2)}+\cdots+p(\alpha)u=G(x)
\tag{B}
\]
and
\[
p(r)=a_0r^n+a_1r^{n-1} + \cdots + a_n
\]
is the characteristic polynomial of the complementary equation
\[
a_0y^{(n)}+a_1y^{(n-1)}+\cdots+a_ny=0.
\]

\item\label{exer:9.3.76}Prove:
\begin{alist}
\item % (a)
The equation
\[
\begin{multlined}
a_0u^{(n)}+\dfrac{p^{(n-1)}(\alpha)}{(n-1)!}u^{(n-1)}+
\dfrac{p^{(n-2)}(\alpha)}{(n-2)!}u^{(n-2)}+\cdots+p(\alpha)u\\
=\left(p_0+p_1x+\cdots+p_kx^k\right)\cos\omega x
+ \left(q_0+q_1x+\cdots+q_kx^k\right)\sin\omega x
\end{multlined}
\tag{A}
\]
has a particular solution of the form
\[
u_p=x^m\left(u_0+u_1x+\cdots+u_kx^k\right)\cos\omega x+
\left(v_0+v_1x+\cdots+v_kx^k\right)\sin\omega x.
\]
\item % (b)
If $\lambda+i\omega$ is a zero of $p$ with multiplicity $m\ge1$, then
(A) can be written as
\[
a(u''+\omega^2 u)=
\left(p_0+p_1x+\cdots+p_kx^k\right)\cos\omega x+
\left(q_0+q_1x+\cdots+q_kx^k\right)\sin\omega x,
\]
which has a particular solution of the form
\[
u_p=U(x)\cos\omega x+V(x)\sin\omega x,
\]
where
\[
U(x)=u_0x+u_1x^2+\cdots+u_kx^{k+1},\,V(x)=v_0x+v_1x^2+\cdots+v_kx^{k+1}
\]
and
\begin{align*}
a(U''(x)+2\omega V'(x))&=p_0+p_1x+\cdots+p_kx^k\\
a(V''(x)-2\omega U'(x))&=q_0+q_1x+\cdots+q_kx^k.
\end{align*}
\end{alist}

\end{exerciselist}


\section{Variation of Parameters for Higher Order  Equations}
\label{sec:higher_variation_param}

\boxit{Derivation of the method}

We assume throughout this section that the nonhomogeneous linear
equation
\begin{equation} \label{eq:9.4.1}
P_0(x)y^{(n)}+P_1(x)y^{(n-1)}+\cdots+P_n(x)y=F(x)
\end{equation}
is normal on an interval $(a,b)$. We'll abbreviate this equation as
$Ly=F$, where
\[
Ly=P_0(x)y^{(n)}+P_1(x)y^{(n-1)}+\cdots+P_n(x)y.
\]
When we speak of solutions of this equation and its complementary
equation $Ly=0$, we mean solutions on $(a,b)$. We'll show how to use
the method of variation of parameters to find a particular solution of
$Ly=F$, provided that we know a fundamental set of solutions
$\{y_1,y_2,\dots,y_n\}$ of $Ly=0$.


We seek a particular solution of $Ly=F$ in the form
\begin{equation} \label{eq:9.4.2}
y_p=u_1y_1+u_2y_2+\cdots+u_ny_n
\end{equation}
where $\{y_1,y_2,\dots,y_n\}$ is a known fundamental set of solutions
of the complementary equation
\[
P_0(x)y^{(n)}+P_1(x)y^{(n-1)}+\cdots+P_n(x)y=0
\]
and $u_1$, $u_2$, \dots, $u_n$ are functions to be determined. We begin by
imposing the following $n-1$ conditions on $u_1,u_2,\dots,u_n$:
\begin{equation} \label{eq:9.4.3}
\begin{aligned}
u'_1y_1+u'_2y_2+&\cdots+u'_ny_n=0 \\
u'_1y'_1+u'_2y'_2+&\cdots+u'_ny'_n=0 \\
\phantom{u'_1y^{(n_1)}+u'_2y_2^{(n-1)}}&\vdots \\
u'_1y_1^{(n-2)}+u'_2y^{(n-2)}_2+&\cdots+u'_ny^{(n-2)}_n=0.
\end{aligned}
\end{equation}
These conditions lead to simple formulas for the first $n-1$
derivatives of $y_p$:
\begin{equation} \label{eq:9.4.4}
y^{(r)}_p=u_1y^{(r)}_1+u_2y_2^{(r)}\cdots+u_ny^{(r)}_n,\ 0
\le r \le n-1.
\end{equation}
These formulas are easy to remember, since they look as though we
obtained them by differentiating \eqref{eq:9.4.2} $n-1$ times while
treating $u_1$, $u_2$, \dots, $u_n$ as constants. To see that
\eqref{eq:9.4.3}
implies \eqref{eq:9.4.4}, we first differentiate \eqref{eq:9.4.2} to obtain
\[
y_p'=u_1y_1'+u_2y_2'+\cdots+u_ny_n'+u_1'y_1+u_2'y_2+\cdots+u_n'y_n,
\]
which  reduces to
\[
y_p'=u_1y_1'+u_2y_2'+\cdots+u_ny_n'
\]
because of the first equation in \eqref{eq:9.4.3}. Differentiating this
yields
\[
y_p''=u_1y_1''+u_2y_2''+\cdots+u_ny_n''+u_1'y_1'+u_2'y_2'+\cdots+u_n'y_n',
\]
which reduces to
\[
y_p''=u_1y_1''+u_2y_2''+\cdots+u_ny_n''
\]
because of  the second equation in \eqref{eq:9.4.3}.
Continuing in this way yields \eqref{eq:9.4.4}.

The last equation in \eqref{eq:9.4.4} is
\[
y_p^{(n-1)}=u_1y_1^{(n-1)}+u_2y_2^{(n-1)}+\cdots+u_ny_n^{(n-1)}.
\]
Differentiating this yields
\[
y_p^{(n)}=u_1y_1^{(n)}+u_2y_2^{(n)}+\cdots+u_ny_n^{(n)}+
u_1'y_1^{(n-1)}+u_2'y_2^{(n-1)}+\cdots+u_n'y_n^{(n-1)}.
\]
Substituting this and  \eqref{eq:9.4.4} into \eqref{eq:9.4.1}
yields
\[
u_1Ly_1+u_2Ly_2+\cdots+u_nLy_n+P_0(x)\left(
u_1'y_1^{(n-1)}+u_2'y_2^{(n-1)}+\cdots+u_n'y_n^{(n-1)}\right)=F(x).
\]
Since  $Ly_i=0$ $(1 \le i \le n)$, this reduces to
\[
u_1'y_1^{(n-1)}+u_2'y_2^{(n-1)}+\cdots+u_n'y_n^{(n-1)}=\frac{F(x)}{P_0(x)}.
\]
Combining this equation with \eqref{eq:9.4.3} shows that
\[
y_p=u_1y_1+u_2y_2+\cdots+u_ny_n
\]
is a solution of \eqref{eq:9.4.1} if
\begin{align*}
u'_1y_1+u'_2y_2+&\cdots+u'_ny_n=0 \\
u'_1y'_1+u'_2y'_2+&\cdots+u'_ny'_n=0 \\
\phantom{u'_1y^{(n_1)}+u'_2y_2^{(n-1)}}&\vdots\\
u'_1y_1^{(n-2)}+u'_2y^{(n-2)}_2+&\cdots+u'_ny^{(n-2)}_n=0 \\
u'_1y^{(n-1)}_1+u'_2y^{(n-1)}_2+&\cdots+u'_n y^{(n-1)}_n=F/P_0,
\end{align*}
which can be written in matrix form as
\begin{equation} \label{eq:9.4.5}
\begin{bmatrix}
y_1&y_2&\cdots&y_n \\
y'_1&y'_2&\cdots&y_n'\\
\vdots&\vdots&\ddots&\vdots\\
y_1^{(n-2)}&y_2^{(n-2)}&\cdots&y_n^{(n-2)}\\
y_1^{(n-1)}&y_2^{(n-1)}&\cdots&y_n^{(n-1)}
\end{bmatrix}
\begin{bmatrix}u_1'\\u_2'\\\vdots\\u_{n-1}'\\u_n'\end{bmatrix}=
\begin{bmatrix}0\\0\\ \vdots\\0\\F/ P_0\end{bmatrix}.
\end{equation}
The determinant of this system is the Wronskian $W$ of
the fundamental set of solutions $\{y_1,y_2,\dots,y_n\}$,
which has no zeros on $(a,b)$, by \zcref{thmtype:9.1.4}.
Solving
\eqref{eq:9.4.5} by Cramer's rule yields
\begin{equation}\label{eq:9.4.6}
u'_j=(-1)^{n-j}\frac{FW_j}{P_0W},\quad 1\le j\le n,
\end{equation}
 where $W_j$ is the Wronskian of the set of
functions obtained by deleting $y_j$ from  $\{y_1,y_2,\dots,y_n\}$
 and keeping the remaining functions in the
same order. Equivalently, $W_j$ is the determinant obtained by
deleting the last row and $j$-th column of $W$.

Having obtained $u_1'$, $u_2'$, \dots, $u_n'$, we can integrate  to
obtain $u_1,\,u_2,\dots,u_n$. As in \zcref{sec:variation_param}, we
take the constants
of integration to be zero, and we drop any linear combination of
$\{y_1,y_2,\dots,y_n\}$ that may appear in $y_p$.

\medskip

\begin{remark*}
For efficiency, it's best to compute $W_1$, $W_2$, \dots, $W_n$
first, and then compute $W$ by expanding in cofactors of the last row;
thus,
\[
W=\sum_{j=1}^n(-1)^{n-j}y_j^{(n-1)}W_j.
\]
\end{remark*}


\boxit{Third Order Equations}

If $n=3$, then
\[
W=\begin{vmatrix}
y_1&y_2&y_3 \\
y'_1&y'_2&y'_3 \\
y''_1&y''_2&y''_3
\end{vmatrix}.
\]
Therefore
\[
W_1=\begin{vmatrix}
y_2&y_3 \\
y'_2&y'_3
\end{vmatrix},
\quad W_2=\begin{vmatrix}
y_1&y_3 \\
y'_1&y'_3
\end{vmatrix},
\quad W_3=\begin{vmatrix}
y_1&y_2 \\
y'_1&y'_2
\end{vmatrix},
\]
and  \eqref{eq:9.4.6}  becomes
\begin{equation} \label{eq:9.4.7}
u'_1=\frac{FW_1}{P_0W},\quad u'_2=-\frac{FW_2}{P_0W},\quad
u'_3=\frac{FW_3}{P_0W}.
\end{equation}



\begin{example}\label{example:9.4.1}
 Find a particular solution of
\begin{equation} \label{eq:9.4.8}
xy'''-y''-xy'+y=8x^2e^x,
\end{equation}
given that  $y_1=x$, $y_2=e^x$, and $y_3=e^{-x}$ form a fundamental
set of solutions of the complementary equation. Then find the general
solution of \eqref{eq:9.4.8}.
 \end{example}

\begin{solution}
We seek a particular solution of \eqref{eq:9.4.8} of the form
\[
y_p=u_1x+u_2e^x+u_3e^{-x}.
\]
The Wronskian of  $\{y_1,y_2,y_3\}$ is
\[
W(x)=\begin{vmatrix}
x&e^x&e^{-x} \\ 1&e^x&-e^{-x} \\ 0&e^x&e^{-x}
\end{vmatrix},
\]
so
\begin{align*}
W_1
&=\begin{vmatrix}
e^x&e^{-x}\\ e^x&-e^{-x}
\end{vmatrix}=-2,\\
W_2
&=\begin{vmatrix}
x&e^{-x}\\1&-e^{-x}
\end{vmatrix}=-e^{-x}(x+1),\\
W_3
&=\begin{vmatrix}
x&e^x\\1&e^x
\end{vmatrix}=e^x(x-1).
\end{align*}
Expanding $W$ by cofactors of the last row yields
\[
W=0W_1-e^x W_2+e^{-x}W_3=0(-2)-e^x\left(-e^{-x}(x+1)\right)
+e^{-x}e^x(x-1)=2x.
\]
 Since $F(x)=8x^2e^x$ and $P_0(x)=x$,
\[
\frac{F}{P_0W}=\frac{8x^2e^x}{x\cdot 2x}=4e^x.
\]
Therefore, from \eqref{eq:9.4.7}
\begin{align*}
u'_1&=\phantom{-} 4e^xW_1=\phantom{-}4e^x(-2)=-8e^x,\\
u_2'&=-4e^xW_2=-4e^x\left(-e^{-x}(x+1)\right)=4(x+1),\\
u_3'&=\phantom{-}4e^xW_3=\phantom{-}4e^x\left(e^x(x-1)\right)=4e^{2x}(x-1).
\end{align*}
 Integrating and taking the constants of
integration to be zero yields
\[
u_1=-8e^x,\quad u_2=2(x+1)^2, u_3=e^{2x}(2x-3).
\]
Hence,
\begin{align*}
y_p&=u_1y_1+u_2y_2+u_3y_3\\
&=(-8e^x)x+e^x(2(x+1)^2)+e^{-x}\left(e^{2x}(2x-3)\right)\\
&=e^x(2x^2-2x-1).
\end{align*}
Since $-e^x$ is a solution of the complementary equation, we redefine
\[
y_p=2xe^x(x-1).
\]
 Therefore the general solution of \eqref{eq:9.4.8} is
\[
y=2xe^x(x-1)+c_1x+c_2e^x+c_3e^{-x}.
\]
\end{solution}


\boxit{Fourth Order Equations}

If $n=4$, then
\[
W=\begin{vmatrix}
y_1&y_2&y_3&y_4 \\
y'_1&y'_2&y'_3&y_4' \\
y''_1&y''_2&y''_3&y_4''\\
y'''_1&y'''_2&y'''_3&y_4'''
\end{vmatrix},
\]
Therefore
\[
W_1=\begin{vmatrix}
y_2&y_3&y_4 \\
y'_2&y'_3&y_4'\\
y''_2&y''_3&y_4''
\end{vmatrix},
\quad W_2=\begin{vmatrix}
y_1&y_3&y_4 \\
y'_1&y'_3&y_4'\\
y''_1&y''_3&y_4''
\end{vmatrix},
\quad W_3=\begin{vmatrix}
y_1&y_2&y_4 \\
y'_1&y'_2&y_4'\\
y''_1&y''_2&y_4''
\end{vmatrix},\quad
 W_4=\begin{vmatrix}
y_1&y_2&y_3 \\
y_1'&y'_2&y_3'\\
y_1''&y''_2&y_3''
\end{vmatrix},
\]
and  \eqref{eq:9.4.6}  becomes
\begin{equation} \label{eq:9.4.9}
u'_1=-\frac{FW_1}{P_0W},\quad u'_2=\frac{FW_2}{P_0W},\quad
u'_3=-\frac{FW_3}{P_0W},\quad u'_4=\frac{FW_4}{P_0W}.
\end{equation}




\begin{example}\label{example:9.4.2}
 Find a particular solution of
\begin{equation} \label{eq:9.4.10}
x^4y^{(4)}+6x^3y'''+2x^2y''-4xy'+4y=12x^2,
\end{equation}
given that  $y_1=x$, $y_2=x^2$, $y_3=1/x$ and $y_4=1/x^2$ form a
fundamental set of solutions of the complementary equation. Then find
the general solution of \eqref{eq:9.4.10} on $(-\infty,0)$ and
$(0,\infty)$.
 \end{example}

\begin{solution}
We seek a particular solution of \eqref{eq:9.4.10} of the form
\[
y_p=u_1x+u_2x^2+\frac{u_3}{x}+\frac{u_4}{x^2}.
\]
The Wronskian of  $\{y_1,y_2,y_3,y_4\}$ is
\[
W(x)=\begin{vmatrix}
x&x^2&1/x&-1/x^2 \\
1&2x&-1/x^2&-2/x^3 \\
0&2&2/x^3&6/x^4\\
0&0&-6/x^4&-24/x^5
\end{vmatrix},
\]
so
\begin{align*}
W_1
&=\begin{vmatrix}
x^2&1/x&1/x^2\\[5pt]2x&-1/x^2&-2/x^3\\2&2/x^3&6/x^4
\end{vmatrix}=-\frac{12}{x^4},&%\\
W_2
&=\begin{vmatrix}
x&1/x&1/x^2\\1&-1/x^2&-2/x^3\\0&2/x^3&6/x^4
\end{vmatrix}=-\frac{6}{x^5},\\
W_3
&=\begin{vmatrix}
x&x^2&1/x^2\\1&2x&-2/x^3\\0&2&6/x^4
\end{vmatrix}=\frac{12}{x^2},  &%\\
W_4
&=\begin{vmatrix}
x&x^2&1/x\\1&2x&-1/x^2\\0&2&2/x^3
\end{vmatrix}=\frac{6}{x}.
\end{align*}
Expanding $W$ by cofactors of the last row yields
\begin{align*}
W
&=-0W_1+0 W_2-\left(-\frac{6}{x^4}\right)W_3+\left(-\frac{24}{x^5}\right)W_4\\
&=\frac{6}{x^4}\frac{12}{x^2}-\frac{24}{x^5}\frac{6}{x}=-\frac{72}{x^6}.
\end{align*}
 Since $F(x)=12x^2$ and $P_0(x)=x^4$,
\[
\frac{F}{P_0W}=\frac{12x^2}{x^4}\left(-\frac{x^6}{72}\right)=-\frac{x^4}{6}.
\]
Therefore, from \eqref{eq:9.4.9},
\begin{align*}
u'_1
&=-\left(-\frac{x^4}{6}\right)W_1=\frac{x^4}{6}\left(-\frac{12}{x^4}\right)=-2,\\
u_2'
&=\phantom{-}-\frac{x^4}{6}W_2=-\frac{x^4}{6}\left(-\frac{6}{x^5}\right)
=\frac{1}{x},\\
u_3'
&=-\left(-\frac{x^4}{6}\right)W_3=\frac{x^4}{6}\frac{12}{x^2}=2x^2,\\
u_4'&=\phantom{-}-\frac{x^4}{6}W_4=-\frac{x^4}{6}\frac{6}{x}=-x^3.
\end{align*}
 Integrating these and taking the constants of
integration to be zero yields
\[
u_1=-2x,\quad u_2=\ln|x|,\quad u_3=\frac{2x^3}{3},
u_4=-\frac{x^4}{4}.
\]
Hence,
\begin{align*}
y_p
&=u_1y_1+u_2y_2+u_3y_3+u_4y_4\\
&=(-2x)x+(\ln|x|)x^2+\frac{2x^3}{3}\frac{1}{x}+\left(-\frac{x^4}{4}\right)
\frac{1}{x^2} \\
&=x^2\ln|x|-\frac{19x^2}{12}.
\end{align*}
Since $-19x^2/12$ is a solution of the complementary equation, we redefine
\[
y_p=x^2\ln|x|.
\]
 Therefore
\[
y=x^2\ln|x|+c_1x+c_2x^2+\frac{c_3}{x}+\frac{c_4}{x^2}
\]
is the general solution of \eqref{eq:9.4.10} on $(-\infty,0)$ and $(0,\infty)$.
\end{solution}

\exercises
In \zcref[range]{exer:9.4.1,exer:9.4.21} find a particular
solution,
given the fundamental set of solutions of the complementary equation.

\begin{exerciselist}
\item\label{exer:9.4.1}$x^3y'''-x^2(x+3)y''+2x(x+3)y'-2(x+3)y=-4x^4$;\quad
$\{x,\,x^2,\,xe^x\}$

\item\label{exer:9.4.2}$y'''+6xy''+(6+12x^2)y'+(12x+8x^3)y=x^{1/2}e^{-x^2}$;\quad
 $\{e^{-x^2},\,xe^{-x^2},\,x^2e^{-x^2}\}$

\item\label{exer:9.4.3}$x^3y'''-3x^2y''+6xy'-6y=2x$;
 \quad  $\{x,x^2,x^3\}$

\item\label{exer:9.4.4}$x^2y'''+2xy''-(x^2+2)y'=2x^2$;\quad   $\{1,\,e^x/x,\,e^{-x}/x\}$

\item\label{exer:9.4.5}$x^3y'''-3x^2(x+1)y''+3x(x^2+2x+2)y'-(x^3+3x^2+6x+6)y=x^4e^{-3x}$;\quad$\{xe^x,\,x^2e^x,\,x^3e^x\}$

\item\label{exer:9.4.6}$x(x^2-2)y'''+(x^2-6)y''+x(2-x^2)y'+(6-x^2)y=2(x^2-2)^2$; \quad
 $\{e^x,\,e^{-x},\,1/x\}$

\item\label{exer:9.4.7}$xy'''-(x-3)y''-(x+2)y'+(x-1)y=-4e^{-x}$; \quad
$\{e^x,\,e^x/x,\,e^{-x}/x\}$

\item\label{exer:9.4.8}$4x^3y'''+4x^2y''-5xy'+2y=30x^2$; \quad  $\{\sqrt x,\,1/\sqrt x,\,x^2\}$

\item\label{exer:9.4.9}$x(x^2-1)y'''+(5x^2+1)y''+2xy'-2y=12x^2$; \quad
$\{x,\,1/(x-1),\,1/(x+1)\}$

\item\label{exer:9.4.10}$x(1-x)y'''+(x^2-3x+3)y''+xy'-y=2(x-1)^2$; \quad
$\{x,\,1/x,e^x/x\}$

\item\label{exer:9.4.11}$x^3y'''+x^2y''-2xy'+2y=x^2$;  \quad  $\{x,\,x^2,\,1/x\}$

\item\label{exer:9.4.12}$xy'''-y''-xy'+y=x^2$; \quad
$\{x,\,e^x,\,e^{-x}\}$

\item\label{exer:9.4.13}$xy^{(4)}+4y'''=6 \ln |x|$; \quad   $\{1,\,x,\,x^2,\,1/x\}$

\item\label{exer:9.4.14}$16x^4y^{(4)}+96x^3y'''+72x^2y''-24xy'+9y=96x^{5/2}$; \quad
$\{\sqrt x,\,1/\sqrt x,\,x^{3/2},\,x^{-3/2}\}$

\item\label{exer:9.4.15}$x(x^2-6)y^{(4)}+2(x^2-12)y'''+x(6-x^2)y''+2(12-x^2)y'=2(x^2-6)^2$; \quad $\{1,\,1/x,\,e^x,\,e^{-x}\}$

\item\label{exer:9.4.16}$x^4y^{(4)}-4x^3y'''+12x^2y''-24xy'+24y=x^4$; \quad
$\{x,\,x^2,\,x^3,\,x^4\}$

\item\label{exer:9.4.17}$x^4y^{(4)}-4x^3y'''+2x^2(6-x^2)y''+4x(x^2-6)y'+(x^4-4x^2+24)y=4x^5e^x$;
\quad $\{xe^x,\,x^2e^x,\,xe^{-x},\,x^2e^{-x}\}$

\item\label{exer:9.4.18}$x^4y^{(4)}+6x^3y'''+2x^2y''-4xy'+4y=12x^2$;   \quad
$\{x,x^2,1/x,1/x^2\}$

\item\label{exer:9.4.19}$xy^{(4)}+4y'''-2xy''-4y'+xy=4e^x$; \quad
$\{e^x,\,e^{-x},\,e^x/x,\,e^{-x}/x\}$

\item\label{exer:9.4.20}$xy^{(4)}+(4-6x)y'''+(13x-18)y''+(26-12x)y'+(4x-12)y=3e^x$; \quad
 $\{e^x,\,e^{2x},\,e^x/x,\,e^{2x}/x\}$

\item\label{exer:9.4.21}$x^4y^{(4)}-4x^3y'''+x^2(12-x^2)y''+2x(x^2-12)y'+2(12-x^2)y=2x^5$;
 $\{x,\,x^2,\,xe^x,\,xe^{-x}\}$

\exercisetext{In \zcref[range]{exer:9.4.22,exer:9.4.33} solve the
initial value problem, given the fundamental set of solutions of the
complementary equation. Where indicated by \CGex, graph the solution.}

\item\label{exer:9.4.22}\CGex
$x^3y'''-2x^2y''+3xy'-3y=4x, \quad  y(1)=4,\quad y'(1)=4,
\quad y''(1)=2$; \quad     $\{x,\,x^3,\,x \ln x\}$

\item\label{exer:9.4.23}$x^3y'''-5x^2y''+14xy'-18y=x^3, \quad  y(1)=0,\quad y'(1)=1,\quad
y''(1)=7$;
\quad $\{x^2,\, x^3,\,x^3 \ln x\}$

\item\label{exer:9.4.24}$(5-6x)y'''+(12x-4)y''+(6x-23)y'+(22-12x)y=-(6x-5)^2e^x$; $y(0)=-4$, \quad  $y'(0)=-\dfrac{3}{2}$,\quad $y''(0)=-19$;\quad
 $\{e^x,\,e^{2x},\,xe^{-x} \}$

\item\label{exer:9.4.25}$x^3y'''-6x^2y''+16xy'-16y=9x^4, \quad  y(1)=2,\quad y'(1)=1,\quad y''(1)=5$; $\{x,\,x^4,\,x^4 \ln |x|\}$

\item\label{exer:9.4.26}\CGex
$(x^2-2x+2)y'''-x^2y''+2xy'-2y=(x^2-2x+2)^2, \quad  y(0)=0,\quad y'(0)=5$,
$y''(0)=0$;\quad       $\{x,\,x^2,\,e^x\}$

\item\label{exer:9.4.27}$x^3y'''+x^2y''-2xy'+2y=x(x+1), \quad  y(-1)=-6,\quad y'(-1)=\dfrac{43}{6},\quad y''(-1)=-\dfrac{5}{2}$;\quad $\{x,\,x^2,\,1/x\}$

\item\label{exer:9.4.28}$(3x-1)y'''-(12x-1)y''+9(x+1)y'-9y=2e^x(3x-1)^2, \quad  y(0)=\dfrac{3}{4}$,
$y'(0)=\dfrac{5}{4}$, \quad  $y''(0)=\dfrac{1}{4}$;\quad
$\{x+1,\,e^x,\,e^{3x}\}$

\item\label{exer:9.4.29}\CGex
$(x^2-2)y'''-2xy''+(2-x^2)y'+2xy=2(x^2-2)^2, \quad  y(0)=1,\quad y'(0)=-5$,
 $y''(0)=5$;\quad      $\{x^2,\,e^x,\,e^{-x}\}$

\item\label{exer:9.4.30}\CGex
$x^4y^{(4)}+3x^3y'''-x^2y''+2xy'-2y=9x^2, \quad  y(1)=-7,\quad y'(1)=
-11,\quad y''(1)=-5$,\;
$y'''(1)=6; \quad \{x,\,x^2,\,1/x,\,x\ln x\}$

\item\label{exer:9.4.31}$(2x-1)y^{(4)}-4xy'''+(5-2x)y''+4xy'-4y=6(2x-1)^2$, \quad  $y(0)=\dfrac{55}{4}$,\quad $y'(0)=0$, $y''(0)=13$, \quad  $y'''(0)=1$;\quad    $\{x,\,e^x,\,e^{-x},\,e^{2x}\}$

\item\label{exer:9.4.32}$4x^4y^{(4)}+24x^3y'''+23x^2y''-xy'+y=6x$,\quad $y(1)=2,\quad
y'(1)=0,\quad y''(1)=4,\quad y'''(1)=-\dfrac{37}{4}$;
\quad $\{x,\sqrt x,1/x,1/\sqrt x\}$

\item\label{exer:9.4.33}$x^4y^{{(4)}}+5x^3y'''-3x^2y''-6xy'+6y=40x^3,
\quad  y(-1)=-1, \;  y'(-1)=-7$,
$y''(-1)=-1,\quad y'''(-1)=-31$;\quad
$\{x,\, x^3,\,1/x,\,1/x^2\}$

\item\label{exer:9.4.34}Suppose the equation
\[
P_0(x)y^{(n)}+P_1(x)y^{(n-1)}+\cdots+P_n(x)y=F(x)
\tag{A}
\]
is normal on an interval $(a,b)$. Let $\{y_1,y_2,\dots,y_n\}$ be a
fundamental set of solutions of its complementary equation on $(a,b)$,
let $W$ be the Wronskian of $\{y_1,y_2,\dots,y_n\}$, and let $W_j$ be
the determinant obtained by deleting the last row and the $j$-th
column of $W$. Suppose $x_0$ is in $(a,b)$, let
\[
u_j(x)=(-1)^{(n-j)}\int_{x_0}^x\frac{F(t)W_j(t)}{P_0(t)W(t)}\,dt,
\quad 1\le j\le n,
\]
and define
\[
y_p=u_1y_1+u_2y_2+\cdots+u_ny_n.
\]
\begin{alist}
\item % (a)
Show that $y_p$ is a solution of (A) and that
\[
y_p^{(r)}=u_1y^{(r)}_1+u_2y_2^{(r)}\cdots+u_ny^{(r)}_n,\quad 1
\le r \le n-1,
\]
and
\[
y_p^{(n)}=u_1y_1^{(n)}+u_2y_2^{(n)}+\cdots+u_ny_n^{(n)}+\frac{F}{P_0}.
\]
\hint{See the derivation of the method of variation of
parameters at the beginning of the section.}
\item % (b)
Show that    $y_p$  is the
solution of the initial value problem
\begin{gather*}
P_0(x)y^{(n)}+P_1(x)y^{(n-1)}+\cdots+P_n(x)y=F(x),\\
y(x_0)=0,\;
y'(x_0)=0,\dots,\; y^{(n-1)}(x_0)=0.
\end{gather*}
\item % (c)
Show that  $y_p$  can be written as
\[
y_p(x)=\int_{x_0}^x G(x,t)F(t)\,dt,
\]
where
\[
G(x,t)=\frac{1}{P_0(t)W(t)}\begin{vmatrix}
y_1(t)&y_2(t)&\cdots&y_n(t)\\
y_1'(t)&y_2'(t)&\cdots&y_n'(t)\\
\vdots&\vdots&\ddots&\vdots\\
y_1^{(n-2)}(t)&y_2^{(n-2)}(t)&\cdots&y_n^{(n-2)}(t)\\
y_1(x)&y_2(x)&\cdots&y_n(x)\end{vmatrix},
\]
which is called  the
\personHref{Green}{Green's function} for (A).
\item % (c)
Show that
\[
\frac{\partial^{j}G(x,t)}{\partial x^j}
=\frac{1}{P_0(t)W(t)}\begin{vmatrix}
y_1(t)&y_2(t)&\cdots&y_n(t)\\
y_1'(t)&y_2'(t)&\cdots&y_n'(t)\\
\vdots&\vdots&\ddots&\vdots\\
y_1^{(n-2)}(t)&y_2^{(n-2)}(t)&\cdots&y_n^{(n-2)}(t)\\
y_1^{(j)}(x)&y_2^{(j)}(x)&\cdots&y_n^{(j)}(x)\end{vmatrix},\quad 0\le j\le
n.
\]
\item % (d)
Show that if $a<t<b$ then
\[
\left.\frac{\partial^{j}G(x,t)}{\partial x^j}\right|_{x=t}=
\begin{cases}
0,&1\le j\le n-2,\\
\dfrac{1}{P_0(t)},&j=n-1.
\end{cases}
\]
\item % (e)
Show that
\[
y_p^{(j)}(x)=\begin{cases}
\dst\int_{x_0}^x\dfrac{\partial^{j}G(x,t)}{\partial x^j}F(t)\,dt,
&1\le j\le n-1,\\
\dst\dfrac{F(x)}{P_0(x)}+
\int_{x_0}^x\frac{\partial^{(n)}G(x,t)}{\partial x^n}F(t)\,dt,
&j=n.
\end{cases}
\]
\end{alist}

\exercisetext{In \zcref[range]{exer:9.4.35,exer:9.4.42}
use the method suggested by \zcref{exer:9.4.34} to find a particular
 solution in the form
$y_p=\int_{x_0}^x G(x,t)F(t)\,dt$, given the indicated fundamental
set of
 solutions.
 Assume that $x$ and $x_0$ are in an interval on which
the equation is normal.}


\item\label{exer:9.4.35}$y'''+2y'-y'-2y=F(x); \quad \{e^x,\,e^{-x},e^{-2x}\}$

\item\label{exer:9.4.36}$x^3y'''+x^2y''-2xy'+2y=F(x); \quad \{x,\,x^2,\,1/x\}$

\item\label{exer:9.4.37}$x^3y'''-x^2(x+3)y''+2x(x+3)y'-2(x+3)y=F(x); \{x,x^2,xe^x\}$

\item\label{exer:9.4.38}$x(1-x)y'''+(x^2-3x+3)y''+xy'-y=F(x); \quad \{x,\,1/x,\,e^x/x\}$

\item\label{exer:9.4.39}$y^{(4)}-5y''+4y=F(x); \quad \{e^x,\,e^{-x},\,e^{2x},\,e^{-2x}\}$

\item\label{exer:9.4.40}$xy^{(4)}+4y'''=F(x); \quad  \{1,\,x,\,x^2,\,1/x\}$

\item\label{exer:9.4.41}$x^4y^{(4)}+6x^3y'''+2x^2y''-4xy'+4y=F(x)$;\quad
$\{x,x^2,1/x,1/x^2\}$

\item\label{exer:9.4.42}$xy^{(4)}-y'''-4xy'+4y'=F(x); \quad \{1,\,x^2,\,e^{2x}, e^{-2x}\}$

\end{exerciselist}
